% !TeX root = 20260911.tex
% !LW recipe=lablog XeLaTeX
\input{_support/style.tex}

\begin{document}

\reporttitle{2026-09-11}{실험 결과 및 설계 기록}

\section*{목적과 이 보고서의 범위}

이 날짜의 실험 설정, 관찰 결과, 원인에 대한 해석과 후속 판단을 함께 정리한다. 다른 날짜의 평가 방식이나 표본 수가 다른 점수를 직접 합치지 않으며, 실행 완료와 논문 근거 채택을 구별한다.

\recordnote{수합·문서 개편: 2026-10-20. 실행 날짜는 각 절과 원시 실행 목록을 따른다. 아래의 과거 계획·진행 상태는 당시 기록이며, 현재 채택 상태와 최종 결과는 후속 실험을 함께 확인한다.}

\section{2026-09-11 · 0911 — 다른 데이터셋으로 옮기면 되는가: NF-v3 4종 데이터 품질 감사 (실험 전 검토)}
\label{report:20260911_0911}
\recordnote{이 절은 당시 설정의 결과다. 평가 방식 변경·정정·후속 비교는 뒤에 명시한 실험과 구별한다.}
\textbf{09-22 정의 정정:} 이 로그의 과거 “현실적 oracle/\allowbreak{}천장” 표현은 말미의 정정을 따른다. Bank 정답 선택, bank+\allowbreak{}다수 라벨 정책, bank 없는 test 라벨 대입, 클래스별 F1 상한은 서로 다른 값이다.



\textbf{결론: bot\_\allowbreak{}iot는 부적합, ton\_\allowbreak{}iot는 조건부(꼬리 2개 클래스·미등장 family 프로토콜만), unsw\_\allowbreak{}nb15는 보조용.} cic2018의 두 데이터 문제(동일 벡터·상반 라벨, 단일 벡터 시나리오)는 ton\_\allowbreak{}iot에서 2배 규모로 반복되고, bot\_\allowbreak{}iot는 공격이 자명하게 분리돼 배울 것이 없으며, unsw\_\allowbreak{}nb15는 benign/\allowbreak{}attack이 TTL 지름길로 갈리고 남은 오차는 IXIA 공격 카테고리 간 혼동이다. 원자료·코드: \texttt{docs/\allowbreak{}\allowbreak{}res\allowbreak{}earch/\allowbreak{}\allowbreak{}20\allowbreak{}260911/\allowbreak{}\allowbreak{}d\allowbreak{}ataset\_\allowbreak{}\allowbreak{}q\allowbreak{}uality\_\allowbreak{}\allowbreak{}a\allowbreak{}udit/\allowbreak{}\allowbreak{}}.

\subsection{1. 전 행 hash 감사 (uncapped suite, 66,935,021행)}
\begin{landscape}
\begingroup
\scriptsize
\begin{longtable}{>{\raggedright\arraybackslash}p{\dimexpr 0.12943\linewidth-2\tabcolsep\relax}>{\raggedleft\arraybackslash}p{\dimexpr 0.10187\linewidth-2\tabcolsep\relax}>{\raggedleft\arraybackslash}p{\dimexpr 0.09793\linewidth-2\tabcolsep\relax}>{\raggedleft\arraybackslash}p{\dimexpr 0.09006\linewidth-2\tabcolsep\relax}>{\raggedleft\arraybackslash}p{\dimexpr 0.11762\linewidth-2\tabcolsep\relax}>{\raggedleft\arraybackslash}p{\dimexpr 0.15305\linewidth-2\tabcolsep\relax}>{\raggedleft\arraybackslash}p{\dimexpr 0.17667\linewidth-2\tabcolsep\relax}>{\raggedleft\arraybackslash}p{\dimexpr 0.13337\linewidth-2\tabcolsep\relax}}
\toprule
데이터셋 & 행 & 고유 벡터 & 중복률 & 상반 라벨 그룹에 속한 행 & test F1 상한(클래스별 최대의 평균) & chrono test가 train에 그대로 있는 비율 & 그중 train 다수 라벨과 충돌 \\ \midrule
\endfirsthead
데이터셋 & 행 & 고유 벡터 & 중복률 & 상반 라벨 그룹에 속한 행 & test F1 상한(클래스별 최대의 평균) & chrono test가 train에 그대로 있는 비율 & 그중 train 다수 라벨과 충돌 \\ \midrule
\endhead
{}bot\_\allowbreak{}iot & 16,933,808 & 669,506 & 0.960 & 6.1\% & 0.966 & 9.7\% & 0.9\% \\
{}ton\_\allowbreak{}iot & 27,520,260 & 4,087,189 & 0.851 & \textbf{60.7\%} & 0.968 & \textbf{76.9\%} & 7.5\% \\
{}unsw\_\allowbreak{}nb15 & 2,365,424 & 1,061,597 & 0.551 & 1.3\% & 0.959 & 57.4\% & 0.3\% \\
{}cse\_\allowbreak{}cic\_\allowbreak{}ids2018 & 20,115,529 & 8,805,692 & 0.562 & 26.6\% & 0.862 & 49.8\% & 1.2\% \\
\bottomrule
\end{longtable}
\endgroup
\end{landscape}
test F1 상한 =\allowbreak{} 시간순 test 행 안에서 같은 벡터에 같은 답을 준다는 제약만 두고 클래스별로 F1을 최대화한 값(순도순 그룹 선택, \texttt{ceilings\allowbreak{}\_\allowbreak{}\allowbreak{}v2.\allowbreak{}py}). 클래스마다 따로 최대화하므로 평균은 macro-F1의 상계다. \textbf{정정(0911 후반)}: 처음 계산한 "hash-majority 상한"은 train+\allowbreak{}test를 합친 다수 라벨 배정이라 F1 상한이 아니었다(BoT theft 0.40 vs 실제 상한 0.98, ToN scanning 0.05 vs 1.00). 그 오류는 검증 워크플로의 리뷰어 렌즈가 잡았고 \texttt{ceilings\allowbreak{}\_\allowbreak{}\allowbreak{}v2.\allowbreak{}csv}로 교체했다. 두 번째로 구분할 양은 "train 라벨과의 충돌"(test 벡터가 train에 있는데 다수 라벨이 다른 행)로, 함수로는 맞힐 수 있지만 train으로 학습한 모델은 틀리는 시점 간 라벨 이동이다: ToN scanning 98\%, BoT theft 75\%, CIC infiltration 50\%, dos 35\%, UNSW backdoor 40\%.

\subsection{2. 데이터셋별 판정}
\textbf{bot\_\allowbreak{}iot — 부적합.} ddos·dos·recon은 상한 0.99이고 문헌상 단순 모델도 99.9\%다. 어려운 두 클래스는 라벨 충돌: benign(51,989행, 전체의 0.3\%)의 64\%가 recon/\allowbreak{}dos와 동일 벡터, theft(1,615행)의 68\%가 recon과 동일 벡터인데 test 창에서는 theft가 다수(상한 0.98)이고 train에서는 recon이 다수라 chrono test의 75\%가 train 라벨과 충돌. full-XGB ddos F1 0.031은 ddos 마지막 20\% 구간의 벡터가 train에 0\% 존재해 같은 도구의 dos로 분류되는 wave 아티팩트. benign이 0.3\%라 benign FPR 축이 의미가 없다.

\textbf{ton\_\allowbreak{}iot — 조건부.} benign 62\%·ddos 78\%·scanning 97\%·dos 44\%·xss 41\%·password 33\%가 상반 라벨 동일 벡터. scanning은 test 창에서는 순수(상한 1.00)하지만 같은 벡터가 train에서 benign이라(충돌 98\%) 학습 불가(0821 확인 재현). 반면 mitm(6,013행, 고유 벡터 90.6\%, test 상한 0.99)과 ransomware(3,971행, 90.8\%, 0.95)는 깨끗한데 full-XGB F1이 0.162 /\allowbreak{} 0.150(정밀도 0.09 /\allowbreak{} 0.08, 재현율 0.73 /\allowbreak{} 0.94) — IR 2,800\textasciitilde{}4,200:1에서의 정밀도 문제로, 실제 headroom이 있는 유일한 꼬리다. backdoor(203k)·injection(382k)은 꼬리가 아니고 이미 0.998 /\allowbreak{} 0.941. family가 날짜 블록(scanning d4  \ensuremath{\rightarrow}  …  \ensuremath{\rightarrow}  backdoor/\allowbreak{}ransomware d9  \ensuremath{\rightarrow}  mitm d10)이라 "나중에 등장하는 공격 +\allowbreak{} 소량 라벨" 프로토콜이 구조적으로 가능하다(s37/\allowbreak{}s38 시대 future-unseen AUROC 0.95가 이것). 단 headline macro-F1은 충돌 정책이 결정하므로 scanning·benign \ensuremath{\leftrightarrow} ddos는 hash 상한과 함께 별도 보고해야 한다.

\textbf{unsw\_\allowbreak{}nb15 — 보조용.} 상반 라벨 1.3\%로 가장 깨끗하고 benign은 공격과 동일 벡터가 0이다. 그러나 MIN\_\allowbreak{}TTL/\allowbreak{}MAX\_\allowbreak{}TTL 두 feature만으로 benign/\allowbreak{}attack 정확도 0.9999(문헌의 TTL OS-fingerprint 지름길), TTL을 빼도 XGB benign F1 1.000. 10클래스(benign 포함) macro-F1은 XGB 0.611(sqrt 가중 0.619) vs test F1 상한 평균 0.959로 격차가 크지만, 혼동은 고유 벡터 위의 카테고리 간 혼동(backdoor \ensuremath{\rightarrow} exploits 315·generic 268·fuzzers 213·recon 105, F1 0.008; dos \ensuremath{\rightarrow} exploits 504, F1 0.36; analysis \ensuremath{\rightarrow} exploits 145)이고 문헌은 exploits/\allowbreak{}shellcode/\allowbreak{}recon이 "notionally identical"이며 카테고리 간 일반화 평가에 부적합하다고 지적한다. 꼬리 test 표본이 worms 32·analysis 246·shellcode 477·backdoor 932행이라 단일 seed 수치는 잡음이다. 모든 family가 두 수집일에 동시에 존재해 미등장 family·drift 구조가 없다. 파이프라인 loader 미연결(추가는 3줄).

\subsection{3. 실험 전에 정할 것}
\par\noindent \textbullet\enspace 모든 데이터셋에서 headline macro-F1은 전 클래스로 보고하되, 클래스별 test F1 상한·상한 대비 비율·train 라벨 충돌률을 옆에 둔다(꼬리 클래스를 빼지 않는다). test는 암기·충돌·신규 벡터로 3분해해 보고한다.


\par\noindent \textbullet\enspace 꼬리(정밀도) 이야기: ton\_\allowbreak{}iot mitm·ransomware가 1순위, cic2018 web이 2순위(bot은 이미 0.998), unsw 공격 카테고리는 주의 문구와 함께 보조.


\par\noindent \textbullet\enspace "새 공격 등장 +\allowbreak{} k개 라벨로 context refresh": ton\_\allowbreak{}iot(d1–3 benign-only 창이 calibration, d9–10 등장 family가 target)와 cic2018(bot이 마지막 날·고유 벡터 100\%; infiltration은 3분해 보고와 함께만). bot\_\allowbreak{}iot는 sanity check 이상 쓰지 않는다.


\par\noindent \textbullet\enspace 문헌 확인: ToN/\allowbreak{}BoT/\allowbreak{}CIC2018은 공격이 날짜 블록이라(arXiv 2503.04404 Table 4) 파이프라인의 family별 시간순 split은 family 내부의 앞뒤를 자른 것이고, "새 family 등장"은 family 순서로 따로 구성해야 한다. family 내부 변화도 있다: BoT ddos test 벡터는 train에 없고(wave), ToN scanning·BoT theft는 같은 벡터의 라벨이 시점에 따라 바뀐다.


\subsection{4. HTML 보고서 (post\_\allowbreak{}tabpfn 최신 항목)}
\texttt{docs/\allowbreak{}\allowbreak{}lab\allowbreak{}log/\allowbreak{}\allowbreak{}repo\allowbreak{}rts/\allowbreak{}\allowbreak{}2026\allowbreak{}0911.\allowbreak{}tex\allowbreak{}} — 데이터셋별로 (i) 충돌 벡터 상위 3개의 실제 feature 값과 라벨·시나리오·기간, (ii) 클래스별 분해(고유/\allowbreak{}benign 동일/\allowbreak{}다른 공격 동일, hash 상한, chrono test 암기·충돌률), (iii) 달성 F1 vs hash-majority 상한(row oracle 신기루 구간 표시), (iv) 시간순 60등분 클래스 구성과 클래스별 시간 범위(train/\allowbreak{}val/\allowbreak{}test), 그리고 랜덤 vs 시간순 split의 조회표 정확도(랜덤 84–98\%, BoT-IoT 98.4\%)와 twin-aware 평가 프로토콜 6항을 담았다. 생성 코드 \texttt{scripts/\allowbreak{}\allowbreak{}common/\allowbreak{}\allowbreak{}r\allowbreak{}esearch\_\allowbreak{}\allowbreak{}analysis\allowbreak{}/\allowbreak{}\allowbreak{}2026091\allowbreak{}1/\allowbreak{}\allowbreak{}build\_\allowbreak{}\allowbreak{}dataset\_\allowbreak{}\allowbreak{}report.\allowbreak{}p\allowbreak{}y}(입력 \texttt{report\_\allowbreak{}\allowbreak{}t\allowbreak{}wins.\allowbreak{}jso\allowbreak{}n}, \texttt{report\_\allowbreak{}\allowbreak{}t\allowbreak{}ime.\allowbreak{}json\allowbreak{}}). post\_\allowbreak{}tabpfn.html 통합본에 09-11 항목·상태 셀·핵심 발견 2건·다음 단계 1건을 추가하고 \texttt{sync\_\allowbreak{}\allowbreak{}htm\allowbreak{}l\_\allowbreak{}\allowbreak{}report\allowbreak{}s.\allowbreak{}py}로 동기화했다.

\subsection{5. 중복을 지우면 오르는가 — EXP40(train만) · EXP41(데이터셋 전체)}
보고서 구성을 "문제(라벨 조합별 충돌 행·\%)  \ensuremath{\rightarrow}  oracle처럼 라우팅 불가  \ensuremath{\rightarrow}  현실적 oracle(같은 bank +\allowbreak{} 충돌 벡터는 다수 라벨)  \ensuremath{\rightarrow}  중복 제거 실험  \ensuremath{\rightarrow}  시간순 구성  \ensuremath{\rightarrow}  판정"으로 재편했다. CIC-IDS2018의 oracle 정의를 bank에 맞춰 정정: 달성 0.762  \ensuremath{\rightarrow}  현실적 oracle 0.804  \ensuremath{\rightarrow}  row-level 0.849 (\texttt{cic\_\allowbreak{}\allowbreak{}bank\allowbreak{}\_\allowbreak{}\allowbreak{}realist\allowbreak{}ic\_\allowbreak{}\allowbreak{}oracl\allowbreak{}e.\allowbreak{}py}; infiltration 0.400  \ensuremath{\rightarrow}  0.666  \ensuremath{\rightarrow}  0.970). 이전 판의 0.856은 bank 없이 정답을 다 아는 경우라 별개 양이다. 0821 §5의 oracle 0.964(exp10 family expert 완벽 배정), 0902의 dense oracle 0.81–0.88(exp29 bank)도 모두 row-level이라 같은 제약을 걸면 내려간다.

\textbf{EXP40 — train만 처리, test는 원본} (\texttt{scripts/\allowbreak{}\allowbreak{}exp40\_\allowbreak{}\allowbreak{}de\allowbreak{}dup\_\allowbreak{}\allowbreak{}xgb.\allowbreak{}py}, run \texttt{results/\allowbreak{}\allowbreak{}20260911\allowbreak{}\_\allowbreak{}\allowbreak{}160256\_\allowbreak{}\allowbreak{}3299866\_\allowbreak{}\allowbreak{}nfv3\_\allowbreak{}\allowbreak{}cic\allowbreak{}2018\_\allowbreak{}\allowbreak{}exp\allowbreak{}40\_\allowbreak{}\allowbreak{}dedup\allowbreak{}\_\allowbreak{}\allowbreak{}xgb}; TabPFN arm \texttt{tabpfn/\allowbreak{}\allowbreak{}s\allowbreak{}cripts/\allowbreak{}\allowbreak{}n\allowbreak{}fv3\_\allowbreak{}\allowbreak{}v3\_\allowbreak{}\allowbreak{}e\allowbreak{}xp40\_\allowbreak{}\allowbreak{}ded\allowbreak{}up\_\allowbreak{}\allowbreak{}tabpf\allowbreak{}n.\allowbreak{}py}, seed 42만 완료 후 EXP41로 교체):

\begingroup
\footnotesize
\begin{longtable}{>{\raggedright\arraybackslash}p{\dimexpr 0.19231\linewidth-2\tabcolsep\relax}>{\raggedleft\arraybackslash}p{\dimexpr 0.14744\linewidth-2\tabcolsep\relax}>{\raggedleft\arraybackslash}p{\dimexpr 0.16026\linewidth-2\tabcolsep\relax}>{\raggedleft\arraybackslash}p{\dimexpr 0.14103\linewidth-2\tabcolsep\relax}>{\raggedleft\arraybackslash}p{\dimexpr 0.14103\linewidth-2\tabcolsep\relax}>{\raggedleft\arraybackslash}p{\dimexpr 0.21795\linewidth-2\tabcolsep\relax}}
\toprule
train 처리 & train 행 & XGB macro-F1 & 암기 벡터 정확도 & 신규 벡터 정확도 & TabPFN C0 macro (s42) \\ \midrule
\endfirsthead
train 처리 & train 행 & XGB macro-F1 & 암기 벡터 정확도 & 신규 벡터 정확도 & TabPFN C0 macro (s42) \\ \midrule
\endhead
{}원본 & 12,069,313 & \textbf{0.7548} & 99.99\% & 99.08\% & 0.7628 \\
{}중복 제거 & 5,263,971 & 0.6493 & 88.97\% & 99.13\% & 0.6000 \\
{}중복 제거 +\allowbreak{} 충돌은 다수 라벨 & 5,250,813 & 0.6914 & 95.91\% & 99.23\% & 0.6563 \\
{}중복 제거 +\allowbreak{} 충돌 벡터 삭제 & 5,237,698 & 0.6910 & 95.92\% & 99.21\% & (log 참조) \\
\bottomrule
\end{longtable}
\endgroup
가설의 조건과 다르지만(test가 원본) 손실 위치를 보여 준다: 신규 벡터 정확도는 불변, train에 있던 벡터의 정확도만 떨어진다. 중복은 tree·retrieval 모델의 sample weight다(EXP16·exp32/\allowbreak{}33과 같은 방향).

\textbf{EXP41 — 데이터셋 전체에서 처리한 뒤 같은 규칙으로 재split} (\texttt{scripts/\allowbreak{}\allowbreak{}exp41\_\allowbreak{}\allowbreak{}de\allowbreak{}dup\_\allowbreak{}\allowbreak{}data\allowbreak{}set\_\allowbreak{}\allowbreak{}xgb.\allowbreak{}py}, run \texttt{results/\allowbreak{}\allowbreak{}20260911\allowbreak{}\_\allowbreak{}\allowbreak{}165*\_\allowbreak{}\allowbreak{}nf\allowbreak{}v3\_\allowbreak{}\allowbreak{}cic20\allowbreak{}18\_\allowbreak{}\allowbreak{}exp41\allowbreak{}\_\allowbreak{}\allowbreak{}dedup\_\allowbreak{}\allowbreak{}d\allowbreak{}ataset\_\allowbreak{}\allowbreak{}x\allowbreak{}gb}; 행 5개 미만 시나리오는 train 전용):

\begin{landscape}
\begingroup
\scriptsize
\begin{longtable}{>{\raggedright\arraybackslash}p{\dimexpr 0.13673\linewidth-2\tabcolsep\relax}>{\raggedleft\arraybackslash}p{\dimexpr 0.10102\linewidth-2\tabcolsep\relax}>{\raggedleft\arraybackslash}p{\dimexpr 0.09592\linewidth-2\tabcolsep\relax}>{\raggedleft\arraybackslash}p{\dimexpr 0.11122\linewidth-2\tabcolsep\relax}>{\raggedleft\arraybackslash}p{\dimexpr 0.11122\linewidth-2\tabcolsep\relax}>{\raggedleft\arraybackslash}p{\dimexpr 0.10102\linewidth-2\tabcolsep\relax}>{\raggedleft\arraybackslash}p{\dimexpr 0.08571\linewidth-2\tabcolsep\relax}>{\raggedleft\arraybackslash}p{\dimexpr 0.08571\linewidth-2\tabcolsep\relax}>{\raggedleft\arraybackslash}p{\dimexpr 0.08571\linewidth-2\tabcolsep\relax}>{\raggedleft\arraybackslash}p{\dimexpr 0.08571\linewidth-2\tabcolsep\relax}}
\toprule
데이터 & train 행 & test 행 & test 내부 충돌 행 & XGB macro-F1 & 현실적 oracle & brute & dos & inf & web \\ \midrule
\endfirsthead
데이터 & train 행 & test 행 & test 내부 충돌 행 & XGB macro-F1 & 현실적 oracle & brute & dos & inf & web \\ \midrule
\endhead
{}원본 & 12,069,313 & 4,023,110 & 788,269 & 0.7548 & 0.856 & 0.916 & 0.782 & 0.429 & 0.172 \\
{}중복 제거 & 5,295,822 & 1,765,277 & 396 & 0.7613 & 0.997 & 0.999 & 1.000 & 0.178 & 0.175 \\
{}중복 제거 +\allowbreak{} 충돌은 다수 라벨 & 5,283,411 & 1,761,143 & 0 & 0.7759 & 1.000 & 0.999 & 1.000 & 0.214 & 0.244 \\
{}중복 제거 +\allowbreak{} 충돌 벡터 삭제 & 5,271,045 & 1,757,021 & 0 & 0.7763 & 1.000 & 0.999 & 1.000 & 0.213 & 0.247 \\
\bottomrule
\end{longtable}
\endgroup
\end{landscape}
판정: macro는 +\allowbreak{}0.02 오르지만 이유는 head다. slowhttptest=\allowbreak{}ftp 단일 벡터 등 라벨 결함 블록이 test에서 사라져 dos·brute·ddos가 1.00이 된다. 꼬리는 반대로 infiltration 0.43  \ensuremath{\rightarrow}  0.18–0.21, web 0.17  \ensuremath{\rightarrow}  0.18–0.25 — 암기로 맞히던 중복 벡터가 빠지고 고유 벡터만 남아 진짜 일반화 난도가 드러난다(0821 EXP18의 inf 0.19·web 0.11과 같은 결과). 현실적 oracle이 1.0이 되므로 남은 격차 0.22는 "데이터가 막는 것"이 아니라 "학습 가능하지만 못 배운 것"이다. 원본 arm의 test는 파이프라인과 4행 차이(재구현 split). TabPFN composed-C0 arm(\texttt{tabpfn/\allowbreak{}\allowbreak{}s\allowbreak{}cripts/\allowbreak{}\allowbreak{}n\allowbreak{}fv3\_\allowbreak{}\allowbreak{}v3\_\allowbreak{}\allowbreak{}e\allowbreak{}xp41\_\allowbreak{}\allowbreak{}ded\allowbreak{}up\_\allowbreak{}\allowbreak{}datas\allowbreak{}et\_\allowbreak{}\allowbreak{}tabpf\allowbreak{}n.\allowbreak{}py}, 4 처리 × seed 42·43·44)은 GPU 체인 16:59 시작. 원본 arm 3 seed는 완료됐으나 dedup 3 arm × 3 seed는 예측을 다 마친 뒤 채점 단계에서 죽었다 — dedup 후에는 "train에 다른 라벨로 있던 벡터" bucket이 비는데 \texttt{precisio\allowbreak{}n\_\allowbreak{}\allowbreak{}recall\allowbreak{}\_\allowbreak{}\allowbreak{}fscore\_\allowbreak{}\allowbreak{}support}가 빈 배열을 거부(\texttt{ValueErr\allowbreak{}or: Foun\allowbreak{}d empty \allowbreak{}input ar\allowbreak{}ray}). 18:36 채점 루프에 빈 bucket guard(요약에 rows=\allowbreak{}0·NaN 기록 후 건너뜀)를 넣고 9 run을 3 lane 병렬(arm당 lane, seed 순차, lane 시작 150 s 간격)로 재실행. 완료된 원본 arm 3 run은 bucket이 모두 비어 있지 않아 guard가 코드 경로에 닿지 않는다(출력 불변). run당 \textasciitilde{}6–7분(test 1.76M행).

\textbf{EXP41 TabPFN 결과} (12 run 완료 19:26; 클래스별 F1은 seed 평균)

\begin{landscape}
\begingroup
\scriptsize
\begin{longtable}{>{\raggedright\arraybackslash}p{\dimexpr 0.12840\linewidth-2\tabcolsep\relax}>{\raggedleft\arraybackslash}p{\dimexpr 0.08759\linewidth-2\tabcolsep\relax}>{\raggedleft\arraybackslash}p{\dimexpr 0.07738\linewidth-2\tabcolsep\relax}>{\raggedleft\arraybackslash}p{\dimexpr 0.07738\linewidth-2\tabcolsep\relax}>{\raggedleft\arraybackslash}p{\dimexpr 0.07738\linewidth-2\tabcolsep\relax}>{\raggedleft\arraybackslash}p{\dimexpr 0.07738\linewidth-2\tabcolsep\relax}>{\raggedleft\arraybackslash}p{\dimexpr 0.07738\linewidth-2\tabcolsep\relax}>{\raggedleft\arraybackslash}p{\dimexpr 0.08759\linewidth-2\tabcolsep\relax}>{\raggedleft\arraybackslash}p{\dimexpr 0.07738\linewidth-2\tabcolsep\relax}>{\raggedleft\arraybackslash}p{\dimexpr 0.07738\linewidth-2\tabcolsep\relax}>{\raggedleft\arraybackslash}p{\dimexpr 0.07738\linewidth-2\tabcolsep\relax}>{\raggedleft\arraybackslash}p{\dimexpr 0.07738\linewidth-2\tabcolsep\relax}}
\toprule
데이터 & test 행 & seed 42 & 43 & 44 & 평균 & tail F1 & XGB on C0 & brute & dos & inf & web \\ \midrule
\endfirsthead
데이터 & test 행 & seed 42 & 43 & 44 & 평균 & tail F1 & XGB on C0 & brute & dos & inf & web \\ \midrule
\endhead
{}원본 & 4,023,110 & 0.7631 & 0.7621 & 0.7609 & \textbf{0.7620} & 0.5498 & 0.7498 & 0.916 & 0.781 & 0.404 & 0.247 \\
{}중복 제거 & 1,765,277 & 0.8062 & 0.7985 & 0.8056 & \textbf{0.8035} & 0.5582 & 0.8110 & 1.000 & 0.997 & 0.163 & 0.529 \\
{}중복 제거 +\allowbreak{} 충돌은 다수 라벨 & 1,761,143 & 0.7994 & 0.8054 & 0.7858 & \textbf{0.7969} & 0.5389 & 0.7965 & 1.000 & 0.998 & 0.163 & 0.471 \\
{}중복 제거 +\allowbreak{} 충돌 벡터 삭제 & 1,757,021 & 0.7975 & 0.7783 & 0.7689 & \textbf{0.7816} & 0.4945 & 0.7916 & 1.000 & 0.998 & 0.194 & 0.308 \\
\bottomrule
\end{longtable}
\endgroup
\end{landscape}
XGB(full train)와 같은 방향이다: 처리하면 macro가 오르지만(0.762  \ensuremath{\rightarrow}  0.78–0.80) 상승분은 head이고, infiltration은 0.40  \ensuremath{\rightarrow}  0.16–0.19로 내려간다. web은 반대로 0.25  \ensuremath{\rightarrow}  0.31–0.53으로 오르는데, 이는 web과 벡터를 공유하던 benign 44,528행이 사라져 정밀도가 회복되기 때문이다(§2 표의 benign/\allowbreak{}web\_\allowbreak{}attacks 쌍). seed 폭은 처리할수록 커진다(원본 0.0023  \ensuremath{\rightarrow}  충돌 벡터 삭제 0.0286): 남은 행이 줄어 C0 추첨의 영향이 커진다. 처리 간 차이(0.78 vs 0.80)는 이 폭과 같은 크기이므로 arm 사이의 순위는 읽지 않는다.

\subsection{09-22 — Oracle 정의 및 expert 평가와의 관계 정정}
CIC EXP39 원본 cache에서 \texttt{row\_\allowbreak{}\allowbreak{}leve\allowbreak{}l\_\allowbreak{}\allowbreak{}oracle\allowbreak{}\_\allowbreak{}\allowbreak{}bank} 0.849159는 global+\allowbreak{}expert 3개의 정답 참조 선택이다. \texttt{realisti\allowbreak{}c\_\allowbreak{}\allowbreak{}oracle\allowbreak{}\_\allowbreak{}\allowbreak{}bank} 0.803623은 충돌 벡터의 test 다수 라벨을 bank에서 고르는 정책 참조값이다. \texttt{realisti\allowbreak{}c\_\allowbreak{}\allowbreak{}oracle\allowbreak{}\_\allowbreak{}\allowbreak{}perfect\allowbreak{}} 0.855907은 bank 없이 비충돌 행에 정답을 직접 대입하고 충돌만 다수 라벨로 답한다. ToN/\allowbreak{}BoT/\allowbreak{}UNSW 감사와 EXP41/\allowbreak{}42/\allowbreak{}44의 같은 이름 참조값도 bank 없는 데이터 진단이다.

다수 라벨 참조값은 macro-F1 상한이 아니다. Bank·동일 벡터 제약에서 클래스별 F1을 각각 최적화한 평균 0.804844는 별도의 느슨한 macro 상계이며 하나의 정책으로 동시 달성된 값도 아니다. Bank 없는 값이 모순 제거 뒤 1.0이 되더라도 train에서 학습·일반화 가능함을 보장하지 않는다. Train 다수 라벨을 반드시 따른다는 제약도 모든 학습기에 적용할 수 없다.

기존의 회수 가능/\allowbreak{}불가능 구간 누적 그림은 서로 다른 정책의 차이를 분해한 것처럼 보이므로, 실제 F1·bank 정답 선택·다수 라벨 참조·클래스별 상한·train 라벨 충돌률을 병렬 표로 교체했다. 정답 참조 선택은 expert 생성 품질의 주 평가에서 제외한다. 0918 수정 보고서와 평가 프로토콜에 담당 배정만 이상화하는 새 oracle 및 식별력 대조를 명시했다. 0911 원본 bank와 0918 모순 제거 bank의 점수는 직접 개선량으로 비교하지 않는다.

\subsection{09-22 EXP56 — 다섯 조건 실행 완료}
CIC EXP47·ToN EXP48의 모순 제거 고정 split과 기존 EXP51 S1V0 정책으로 \textbf{global /\allowbreak{} 새 담당 배정 oracle /\allowbreak{} expert별 전면 적용 /\allowbreak{} scorer-only /\allowbreak{} scorer+\allowbreak{}verifier}를 실행했다. 기존 정답 참조 union은 주 비교에서 제외했다. 새 학습·독립 seed 재실행이 아닌 frozen prediction 평가다.

\begingroup
\footnotesize
\begin{longtable}{>{\raggedright\arraybackslash}p{\dimexpr 0.34591\linewidth-2\tabcolsep\relax}>{\raggedleft\arraybackslash}p{\dimexpr 0.32704\linewidth-2\tabcolsep\relax}>{\raggedleft\arraybackslash}p{\dimexpr 0.32704\linewidth-2\tabcolsep\relax}}
\toprule
Macro-F1 & CIC2018 & ToN \\ \midrule
\endfirsthead
Macro-F1 & CIC2018 & ToN \\ \midrule
\endhead
{}Global & 0.810653 & 0.683734 \\
{}새 담당 배정 oracle & 0.867295 & 0.687397 \\
{}Expert 전면 적용 범위 & 0.700074–0.794586 & 0.620527–0.685941 \\
{}Scorer-only (동일 호출) & 0.810653 & 0.754919 \\
{}Scorer+\allowbreak{}verifier & 0.810653 & 0.755490 \\
\bottomrule
\end{longtable}
\endgroup
ToN에서 V는 교정 183건을 거절하는 대신 훼손 475건을 막았다. CIC는 저장된 calibration 정책이 호출 0을 선택하므로 S-only도 0이며 V 효과를 식별하지 못한다. Class×expert·영역×expert·H/\allowbreak{}D와 확률 AP, validation 고정 FPR 운영점의 test 전이, 경험적 test ROC의 R@FPR 0.1\%/\allowbreak{}1\%를 추가 계산했다. 주 정책의 예측을 test 임계값으로 변경하지 않았다.

ToN scanning e2·ransomware e5는 AP 및 R@FPR 0.1\%에서도 global·타 expert보다 높다. CIC infiltration e6은 담당 recall이 높아도 AP·동일 FPR recall은 global보다 낮다. 따라서 담당 배정 oracle 점수와 expert 식별력을 따로 제시한다. 세 관점의 관측 품질 분석은 이번 결과로 가능하며, 생성 규칙의 원인을 입증하려면 동일 예산의 무작위/\allowbreak{}균형/\allowbreak{}β 대조와 seed 반복·시간/\allowbreak{}그룹별 불확실성이 추가 필요하다.

결과: \texttt{tabpfn/\allowbreak{}\allowbreak{}r\allowbreak{}esults/\allowbreak{}\allowbreak{}e\allowbreak{}xp56\_\allowbreak{}\allowbreak{}fiv\allowbreak{}e\_\allowbreak{}\allowbreak{}condit\allowbreak{}ion\_\allowbreak{}\allowbreak{}2026\allowbreak{}0922\_\allowbreak{}\allowbreak{}fin\allowbreak{}al/\allowbreak{}\allowbreak{}\{\allowbreak{}cic2\allowbreak{}018,\allowbreak{}toni\allowbreak{}ot\}\allowbreak{}}. 전체 표와 한계: EXP56 결과. 재현 스크립트 \texttt{tabpfn/\allowbreak{}\allowbreak{}s\allowbreak{}cripts/\allowbreak{}\allowbreak{}e\allowbreak{}xp56\_\allowbreak{}\allowbreak{}fiv\allowbreak{}e\_\allowbreak{}\allowbreak{}condit\allowbreak{}ion\_\allowbreak{}\allowbreak{}eval\allowbreak{}uation.\allowbreak{}p\allowbreak{}y}, 보고서 생성 \texttt{scripts/\allowbreak{}\allowbreak{}build\_\allowbreak{}\allowbreak{}ex\allowbreak{}p56\_\allowbreak{}\allowbreak{}repo\allowbreak{}rt\_\allowbreak{}\allowbreak{}secti\allowbreak{}on.\allowbreak{}py}. 기존 S+\allowbreak{}V 예측의 행 단위 재현과 과학적 불변조건 단위 검증 5개를 통과했다.

\subsubsection{2026-09-22 HTML 재작성}
다섯 조건과 세 진단 관점을 도입에 배치하고, 전체 bank의 충분성 판단·개선 과제를 명시했다. 클래스 표는 test support 내림차순, expert는 번호 순서로 유지하며 다섯 조건의 표시 정밀도 최고값(동률 포함)을 강조했다. Macro-F1 비교·고정 FPR 식별력 차이·context 구성 그림을 추가하고 Noto Sans CJK KR을 내장했다.

생성 코드를 확인한 결과 K는 클래스 수로 고정한 것이 아니라 2/\allowbreak{}4/\allowbreak{}8 후보의 tune 정답 기반 최소 손실 선택 지표로 결정됐다. CIC는 8개, ToN은 8개 중 최소 tune 영역 표본 조건 미달 1개를 제거해 7개다. 평가용 지배 클래스 매핑과 실제 생성·scorer 정책을 구별하며 K 선택 기준 자체를 개선 과제로 제시한다. 숫자·bank 재학습 없이 EXP56 final CSV를 재사용했다.

\subsubsection{2026-09-22 후속 수정 — 정답 배정 제외, K<C, 지표의 검증 목적}
이전 개정의 정답 클래스 배정은 상수 expert의 높은 점수를 허용하므로 주 평가에서 제외했다. 2번은 정답 없는 고정 centroid 영역 배정(S/\allowbreak{}V 없음)으로 교체: CIC macro-F1 0.747722, ToN 0.661978. 생성 코드에 1<=\allowbreak{}K<C를 추가했고 기존 CIC 8/\allowbreak{}7 bank는 제약 적용 전 진단이라는 점을 표시했다. Bank 재학습은 하지 않았다.

CIC 동일 성능은 calibration 65개 후보의 benign FPR/\allowbreak{}tail F1 제약 위반과 호출 0의 결과다. Infiltration 정답 배정 0.9140은 e6의 TP 7,556을 유지하면서 FP 806,387 대신 다른 예측기의 FP 236만 남긴 결과다. ToN scanning S+\allowbreak{}V는 TP 4,245/\allowbreak{}FP 172/\allowbreak{}FN 1,784로 0.8128이며 모델별 예측 출처를 검증했다.

추가 계산: 학습 route 다수 클래스를 출력하는 영역별 상수 대조. HTML에서 expert 개수의 상세 결정 경위와 라우팅 회수율 절을 제거하고, 전체 test/\allowbreak{}AP/\allowbreak{}동일 FPR recall/\allowbreak{}val 임계값 전이/\allowbreak{}영역 대조의 목적을 설명했다. 필요한 근거를 동일 예산·클래스 구성 context 대조, global 보정 대조, 새 시간 운영점 전이, 독립 학습 seed 반복으로 구체화했다. EXP52–54의 기존 부분 대조는 보존·구분했다. 결과: \texttt{tabpfn/\allowbreak{}\allowbreak{}r\allowbreak{}esults/\allowbreak{}\allowbreak{}e\allowbreak{}xp56\_\allowbreak{}\allowbreak{}lab\allowbreak{}el\_\allowbreak{}\allowbreak{}free\_\allowbreak{}\allowbreak{}20260922\allowbreak{}/\allowbreak{}\allowbreak{}}.

\subsection{2026-09-22 — EXP57 실행 및 EXP58 모델 상한 추가}
현재 2번은 정답 없는 고정 영역 배정이며 oracle이 아니다. 기존 bank에서 가능한 모델 상한은 EXP58의 6번 bank oracle로 별도 측정했다. Accuracy뿐 아니라 모두 틀린 행의 오탐 목적지까지 최적화한 Macro-F1 상한은 CIC 0.991366, ToN 0.786694다. 같은 예측 bank와 test에서의 상한이며 새 K=\allowbreak{}4 bank의 결과가 아니다. 상한 정의·검증.

EXP57은 K=\allowbreak{}4, seed 42·43·44의 실제 context fit과 세 생성 방식 대조를 진행 중이다. Seed 42의 두 로컬 작업이 native segmentation fault로 종료돼 외부 GPU에 배정했고, 로컬은 seed 43·44를 이어간다. 외부 실행 명령. 이번 설명·상한 결과는 사용자 요청에 따라 HTML에 추가하지 않았다.



\section{NF-v3 suite data-quality audit (2026-09-11, read-only)}
분석 질문은 CIC2018에서 확인한 특징이 동일하지만 라벨이 다른 사례, 단일 특징 벡터에 집중된 공격, 시간 분할 사이의 중복 문제가 다른 세 데이터에도 있는지다. NF-v3 통합 입력 66,935,021행·46개 특징에서 NaN/무한대를 처리한 float32 행의 해시로 중복을 확인했다.

Files

\par\noindent \textbullet\enspace \texttt{audit\_\allowbreak{}\allowbreak{}su\allowbreak{}mmary.\allowbreak{}cs\allowbreak{}v} per dataset: rows, unique vectors, duplicate rate, fraction of rows in mixed-label hash groups, hash-majority macro-F1 ceiling, chronological-test (per-family 60/\allowbreak{}20/\allowbreak{}20 by time) memorization and conflict rates.


\par\noindent \textbullet\enspace \texttt{audit\_\allowbreak{}\allowbreak{}cl\allowbreak{}asses.\allowbreak{}cs\allowbreak{}v} per class: twin fractions (any other label /\allowbreak{} benign), hash-majority F1, time quantiles, chrono-test memorization/\allowbreak{}conflict.


\par\noindent \textbullet\enspace \texttt{audit\_\allowbreak{}\allowbreak{}sc\allowbreak{}enarios.\allowbreak{}csv} per scenario: unique vectors, top-vector share, rows sharing vectors with other scenarios.


\par\noindent \textbullet\enspace \texttt{audit\_\allowbreak{}\allowbreak{}ce\allowbreak{}ilings.\allowbreak{}c\allowbreak{}sv} per class: clean-vector share, majority-assignment and rarest-class- assignment F1 ceilings (all rows and chrono-test rows).


\par\noindent \textbullet\enspace \texttt{unsw\_\allowbreak{}\allowbreak{}xgb\allowbreak{}\_\allowbreak{}\allowbreak{}plain.\allowbreak{}c\allowbreak{}sv}, \texttt{unsw\_\allowbreak{}\allowbreak{}xgb\allowbreak{}\_\allowbreak{}\allowbreak{}sqrt\_\allowbreak{}\allowbreak{}ba\allowbreak{}lanced.\allowbreak{}c\allowbreak{}sv}: XGBoost (300 trees, depth 8, lr 0.05, seed 42) on unsw\_\allowbreak{}nb15 per-family chronological 60/\allowbreak{}20/\allowbreak{}20; F1 on clean vectors only; ceilings.


\par\noindent \textbullet\enspace \texttt{unsw\_\allowbreak{}\allowbreak{}pro\allowbreak{}be.\allowbreak{}log}: TTL-only benign/\allowbreak{}attack shortcut check and XGB confusion matrices with and without MIN\_\allowbreak{}TTL/\allowbreak{}MAX\_\allowbreak{}TTL. Scripts are the exact code that produced the CSVs; they write nothing under the repo.


\par\noindent \textbullet\enspace \texttt{report\_\allowbreak{}\allowbreak{}d\allowbreak{}ata\_\allowbreak{}\allowbreak{}twin\allowbreak{}s.\allowbreak{}py}  \ensuremath{\rightarrow}  \texttt{report\_\allowbreak{}\allowbreak{}t\allowbreak{}wins.\allowbreak{}jso\allowbreak{}n}: per dataset the top-10 conflicting hash groups (labels, scenarios, time spans, the full 46-feature vector) and per-class clean /\allowbreak{} benign-twin /\allowbreak{} other-twin decomposition; \texttt{report\_\allowbreak{}\allowbreak{}d\allowbreak{}ata\_\allowbreak{}\allowbreak{}time\allowbreak{}.\allowbreak{}py}  \ensuremath{\rightarrow}  \texttt{report\_\allowbreak{}\allowbreak{}t\allowbreak{}ime.\allowbreak{}json\allowbreak{}}: class share over 60 time-ordered bins, per-class time spans with the per-family 60/\allowbreak{}20/\allowbreak{}20 boundaries, and the hash-lookup classifier under random 80/\allowbreak{}20 vs chronological split.


\par\noindent \textbullet\enspace \texttt{build\_\allowbreak{}\allowbreak{}da\allowbreak{}taset\_\allowbreak{}\allowbreak{}re\allowbreak{}port.\allowbreak{}py}: renders \texttt{docs/\allowbreak{}\allowbreak{}lab\allowbreak{}log/\allowbreak{}\allowbreak{}repo\allowbreak{}rts/\allowbreak{}\allowbreak{}2026\allowbreak{}0911.\allowbreak{}tex\allowbreak{}} from the above.




\section{문헌 링크 모음 — 모순 · 중복 · 누수 (2026-09-14 조사)}
용어는 보고서와 통일한다.

\par\noindent \textbullet\enspace \textbf{중복}: 같은 (feature 벡터, 라벨) 쌍이 2회 이상 나타남


\par\noindent \textbullet\enspace \textbf{모순}: 한 feature 벡터가 서로 다른 라벨을 가짐


\par\noindent \textbullet\enspace \textbf{누수}: test의 벡터가 train에도 존재함 (분할에 따라 달라지는 성질)


각 항목의 "확인" 열은 실제로 페이지나 PDF를 열어 본문을 확인했는지를 뜻한다.


\medskip\hrule\medskip
\subsection{A. 실제로 행을 지운 연구}
\begingroup
\footnotesize
\begin{longtable}{>{\raggedright\arraybackslash}p{\dimexpr 0.22179\linewidth-2\tabcolsep\relax}>{\raggedright\arraybackslash}p{\dimexpr 0.19615\linewidth-2\tabcolsep\relax}>{\raggedright\arraybackslash}p{\dimexpr 0.22179\linewidth-2\tabcolsep\relax}>{\raggedright\arraybackslash}p{\dimexpr 0.22179\linewidth-2\tabcolsep\relax}>{\raggedright\arraybackslash}p{\dimexpr 0.13846\linewidth-2\tabcolsep\relax}}
\toprule
논문 & 연도 /\allowbreak{} venue & 링크 & 무엇을 지웠나 & 확인 \\ \midrule
\endfirsthead
논문 & 연도 /\allowbreak{} venue & 링크 & 무엇을 지웠나 & 확인 \\ \midrule
\endhead
{}Tavallaee, Bagheri, Lu, Ghorbani — \emph{A Detailed Analysis of the KDD CUP 99 Data Set} & 2009, IEEE CISDA & https:/\allowbreak{}/\allowbreak{}www.ee.torontomu.ca/\allowbreak{}\textasciitilde{}bagheri/\allowbreak{}papers/\allowbreak{}cisda.pdf & \textbf{중복}. KDD'99의 반복 레코드를 1개만 남김  \ensuremath{\rightarrow}  NSL-KDD. train 4,898,431  \ensuremath{\rightarrow}  1,074,992 (78.05\%↓), test 311,027  \ensuremath{\rightarrow}  77,289 (75.15\%↓) & 본문 \\
{}Raskovalov, Gabdullin, Dolmatov — \emph{Investigation and rectification of NIDS datasets and standardized feature set derivation...} & 2022, arXiv only & https:/\allowbreak{}/\allowbreak{}arxiv.org/\allowbreak{}abs/\allowbreak{}2212.13994 & \textbf{모순}. "in NF-ToN-IoT a number of flows are duplicated and labeled as corresponding to different attacks". 21,978,630  \ensuremath{\rightarrow}  18,902,360 (14\%↓). \textbf{mitm 1,052  \ensuremath{\rightarrow}  0} & 본문 +\allowbreak{} Table 1 \\
{}Chen, Tran, Thumati, Bhuyan, Ding — \emph{Data Curation and Quality Assurance for ML-based Cyber Intrusion Detection} & 2021, arXiv (ACM JDIQ 심사 중 표기) & https:/\allowbreak{}/\allowbreak{}arxiv.org/\allowbreak{}abs/\allowbreak{}2105.10041 & \textbf{모순}. 두 클래스에 동시에 존재하는 시퀀스 제거. 제거 전후 FPR 최대 27.5배 차이. HIDS라 직접 비교는 아님. "중복 자체는 결함이 아니고 train/\allowbreak{}test 겹침이 결함"이라는 논지도 이 논문 & 본문 \\
{}Mondragon, Branco, Jourdan, Gutierrez-Rodriguez, Biswal — \emph{Advanced IDS: a comparative study of datasets and ML algorithms...} & 2025, Applied Intelligence & https:/\allowbreak{}/\allowbreak{}link.springer.com/\allowbreak{}article/\allowbreak{}10.1007/\allowbreak{}s10489-025-06422-4 & \textbf{중복}. NetFlow 판 행 수만 보고: BoT-IoT 73,370,443  \ensuremath{\rightarrow}  37,643,287 (49\%↓), ToN-IoT 22,339,021  \ensuremath{\rightarrow}  16,940,365 (24\%↓). 클래스별 분해 없음 & 요약 \\
\bottomrule
\end{longtable}
\endgroup
\subsection{B. 지우지 않고 측정하거나 재라벨한 연구}
\begingroup
\footnotesize
\begin{longtable}{>{\raggedright\arraybackslash}p{\dimexpr 0.21950\linewidth-2\tabcolsep\relax}>{\raggedright\arraybackslash}p{\dimexpr 0.21950\linewidth-2\tabcolsep\relax}>{\raggedright\arraybackslash}p{\dimexpr 0.21950\linewidth-2\tabcolsep\relax}>{\raggedright\arraybackslash}p{\dimexpr 0.21950\linewidth-2\tabcolsep\relax}>{\raggedright\arraybackslash}p{\dimexpr 0.12201\linewidth-2\tabcolsep\relax}}
\toprule
논문 & 연도 /\allowbreak{} venue & 링크 & 무엇을 했나 & 확인 \\ \midrule
\endfirsthead
논문 & 연도 /\allowbreak{} venue & 링크 & 무엇을 했나 & 확인 \\ \midrule
\endhead
{}Jerabek, Luxemburk, Plny, Koumar, Pesek, Hynek — \emph{When Simple Model Just Works: Is Network Traffic Classification in Crisis?} & 2025, arXiv & https:/\allowbreak{}/\allowbreak{}arxiv.org/\allowbreak{}abs/\allowbreak{}2506.08655 & \textbf{우리 "현실적 oracle"과 같은 개념의 선행연구.} "most datasets contain over 50\% redundant samples (identical packet sequences), which frequently appear in both training and test sets"; 중복이 "reduce the theoretical maximum accuracy when identical flows have conflicting labels". 12개 데이터셋(UNSW-NB15, CIC-IDS-2017 포함) & 초록 \\
{}Xie, Li, Zhang, Sun, Xu — \emph{Analysis and Detection against Network Attacks in the Overlapping Phenomenon of Behavior Attribute} & Computers \& Security 123 (2022) & https:/\allowbreak{}/\allowbreak{}arxiv.org/\allowbreak{}abs/\allowbreak{}2310.10660 & \textbf{모순을 삭제 대신 multi-label로 재정의.} "multiple samples with the same features but different labels". UNSW-NB15에서 라벨 조합 57종, 샘플당 평균 라벨 1.689개 & 초록 \\
{}Flood, Engelen, Aspinall, Desmet — \emph{Bad Design Smells in Benchmark NIDS Datasets} & 2024, IEEE EuroS\&P (distinguished paper) & https:/\allowbreak{}/\allowbreak{}distrinet.\allowbreak{}cs.\allowbreak{}kuleuven.\allowbreak{}be/\allowbreak{}news/\allowbreak{}bad\_\allowbreak{}smells\_\allowbreak{}euros\_\allowbreak{}p\_\allowbreak{}template-4.pdf & 삭제 없이 \textbf{지표화}. wrong-label(ENN 기반) CIC-17 infiltration 0.81, CIC-18 infiltration 0.65, CTU-13 Murlo 1.00. traffic collapse(유사도 0.95 초과 쌍 비율) ToN-IoT backdoor 1.00, CIC-18 bot 0.99. 권고는 삭제가 아니라 "avoiding using training and test attack data from the same class and dataset" & PDF 전문 \\
{}Al-Daweri, Zainol Ariffin, Abdullah, Md. Senan — \emph{An Analysis of the KDD99 and UNSW-NB15 Datasets for the IDS} & 2020, Symmetry 12(10):1666 & https:/\allowbreak{}/\allowbreak{}www.mdpi.com/\allowbreak{}2073-8994/\allowbreak{}12/\allowbreak{}10/\allowbreak{}1666 & \textbf{중복} 감사만. UNSW-NB15 공식 train의 42.24\%가 중복, test는 0.00\%. Generic 89.54\%, DoS 68.95\% & 요약 \\
{}Engelen, Rimmer, Joosen — \emph{Troubleshooting an Intrusion Detection Dataset: the CICIDS2017 Case Study} & 2021, IEEE SPW (WTMC) & https:/\allowbreak{}/\allowbreak{}intrusion-detection.distrinet-research.be/\allowbreak{}WTMC2021/\allowbreak{}index.html & \textbf{모순의 원인} 규명. IP +\allowbreak{} 시간 창만으로 라벨링해 거절된 연결에도 공격 라벨. CICFlowMeter FIN 버그로 생긴 2-packet 부산물 flow가 전체의 25.9\% & 요약 \\
{}Liu, Engelen, Lynar, Essam, Joosen — \emph{Error Prevalence in NIDS datasets: A Case Study on CIC-IDS-2017 and CSE-CIC-IDS-2018} & 2022, IEEE CNS, pp. 254–262 & https:/\allowbreak{}/\allowbreak{}intrusion-detection.distrinet-research.be/\allowbreak{}CNS2022/\allowbreak{}CSECICIDS2018.html & 라벨링 로직을 역설계·교정해 공개. \textbf{주의}: 흔히 인용되는 오류율 6.67\% /\allowbreak{} 7.53\%는 2차 인용이고 원문 미확인 & 부분 \\
{}Lanvin, Gimenez, Han, Majorczyk, Mé, Totel — \emph{Errors in the CICIDS2017 Dataset and the Significant Differences in Detection Performances It Makes} & CRiSIS 2022, LNCS 13857 & https:/\allowbreak{}/\allowbreak{}pfgimenez.fr/\allowbreak{}publications/\allowbreak{}2022-CRiSIS/\allowbreak{} & packet misorder, \textbf{packet duplication}, 미라벨 공격을 교정하고 성능 차이 측정. \textbf{주의}: 순환 인용되는 "패킷 5\% 중복", "recall 80 \ensuremath{\rightarrow} 100\%"는 미확인 & 초록 \\
{}Pekár, Jozsa — \emph{Evaluating ML-Based Anomaly Detection Across Datasets of Varied Integrity: A Case Study} & 2024, Computer Networks & https:/\allowbreak{}/\allowbreak{}arxiv.org/\allowbreak{}abs/\allowbreak{}2401.16843 & \textbf{우리가 넘어야 할 반증.} CICIDS2017 PCAP에서 중복 패킷(3.6–6.1\%)을 지우고 flow 재생성했으나 Random Forest 성능은 "exceptional robustness" & 요약 \\
{}Cantone, Marrocco, Bria — \emph{On the Cross-Dataset Generalization of Machine Learning for Network Intrusion Detection} & 2024, IEEE Access 12 & https:/\allowbreak{}/\allowbreak{}arxiv.org/\allowbreak{}abs/\allowbreak{}2402.10974 & 높은 내부 성능을 "sample redundancies"로 귀인. \textbf{DoS Slowhttptest는 105,550 샘플에 고유 인스턴스가 56개} & 요약 \\
\bottomrule
\end{longtable}
\endgroup
\subsection{C. "표준 관행이 아니다"의 근거}
\begingroup
\footnotesize
\begin{longtable}{>{\raggedright\arraybackslash}p{\dimexpr 0.25603\linewidth-2\tabcolsep\relax}>{\raggedright\arraybackslash}p{\dimexpr 0.23190\linewidth-2\tabcolsep\relax}>{\raggedright\arraybackslash}p{\dimexpr 0.25603\linewidth-2\tabcolsep\relax}>{\raggedright\arraybackslash}p{\dimexpr 0.25603\linewidth-2\tabcolsep\relax}}
\toprule
논문 & 연도 /\allowbreak{} venue & 링크 & 핵심 문장 \\ \midrule
\endfirsthead
논문 & 연도 /\allowbreak{} venue & 링크 & 핵심 문장 \\ \midrule
\endhead
{}Flood 외 (위와 동일) & 2024, EuroS\&P & (위) & 38편 검토. "No papers in our overview commented on this or amended their evaluation process, and all were seemingly unaware of these discrepancies." \\
{}Apruzzese, Laskov, Schneider — \emph{SoK: Pragmatic Assessment of Machine Learning for Network Intrusion Detection} & 2023, IEEE EuroS\&P & https:/\allowbreak{}/\allowbreak{}arxiv.org/\allowbreak{}abs/\allowbreak{}2305.00550 & 상위 venue 30편 중 "none of the 30 papers considered different preprocessing mechanisms" \\
{}Leevy, Khoshgoftaar — \emph{A survey and analysis of intrusion detection models based on CSE-CIC-IDS2018 Big Data} & 2020, Journal of Big Data 7:104 & https:/\allowbreak{}/\allowbreak{}doi.org/\allowbreak{}10.1186/\allowbreak{}s40537-020-00382-x & "information on the data cleaning of CSE-CIC-IDS2018 was inadequate across the board" \\
{}Goldschmidt, Chudá — \emph{Network Intrusion Datasets: A Survey, Limitations, and Recommendations} & 2025, Computers \& Security 156 & https:/\allowbreak{}/\allowbreak{}arxiv.org/\allowbreak{}abs/\allowbreak{}2502.06688 & 89개 데이터셋 검토. "duplicated features or entries … These issues are frequently unreported". 정작 자신들의 10단계 제작 권고에는 dedup 단계가 없음 \\
{}Ring, Wunderlich, Scheuring, Landes, Hotho — \emph{A Survey of Network-based Intrusion Detection Data Sets} & 2019, Computers \& Security 86 & https:/\allowbreak{}/\allowbreak{}arxiv.org/\allowbreak{}abs/\allowbreak{}1903.02460 & 중복 제거를 KDD 전용 처방으로만 언급. 15개 데이터셋 속성 표에 중복 항목 없음 \\
{}Arp 외 — \emph{Dos and Don'ts of Machine Learning in Computer Security} & 2022, USENIX Security & https:/\allowbreak{}/\allowbreak{}arxiv.org/\allowbreak{}abs/\allowbreak{}2010.09470 & \textbf{우리에게 걸리는 경고.} "selective snooping describes the cleansing of data based on information not available in practice". 행 삭제를 처방하지 않음 \\
\bottomrule
\end{longtable}
\endgroup
\subsection{D. 우리 데이터(NetFlow v2/\allowbreak{}v3)에 대한 선행연구}
\par\noindent \textbullet\enspace Sarhan, Layeghy, Portmann — \emph{Towards a Standard Feature Set for NIDS Datasets}, https:/\allowbreak{}/\allowbreak{}arxiv.org/\allowbreak{}abs/\allowbreak{}2101.11315 : 본문에 duplicate /\allowbreak{} redundant /\allowbreak{} identical /\allowbreak{} conflict 어느 단어도 없음.


\par\noindent \textbullet\enspace Luay, Layeghy, Hosseininoorbin, Sarhan, Moustafa, Portmann — \emph{Temporal Analysis of NetFlow Datasets for NIDS}, https:/\allowbreak{}/\allowbreak{}arxiv.org/\allowbreak{}abs/\allowbreak{}2503.04404 : 시간 feature 추가. 중복·모순·누수 분석 없음.


\par\noindent \textbullet\enspace \textbf{결론: NF-v2에 대한 중복/\allowbreak{}모순 감사는 Raskovalov 외의 한 문장이 전부이고, NF-v3에 대해서는 전무하다.}


\subsection{E. 인용하지 말 것 (확인 결과 무관)}
\par\noindent \textbullet\enspace Sommer \& Paxson, \emph{Outside the Closed World} (IEEE S\&P 2010) — 중복 관련 서술 없음.


\par\noindent \textbullet\enspace Rosay 외 (ICISSP 2022) — 여기서 말하는 duplication은 중복 \textbf{열}(feature)이지 행이 아님.


\par\noindent \textbullet\enspace Zoghi \& Serpen (arXiv 2101.05067) — PCA/\allowbreak{}t-SNE 투영 상의 기하학적 겹침이지 동일 벡터가 아님.


\par\noindent \textbullet\enspace Peterson 외 \emph{A Review and Analysis of the Bot-IoT Dataset} (IEEE SOSE 2021) 및 \emph{Composition analysis of the Bot-IoT dataset} (IJITCA 2022) — 유료라 미확인. 인용 금지.




\section{추가 분석과 수치 검토}
이 보고서는 09-11\textasciitilde{}15의 데이터 감사와 당시 권고를 보존한 기록입니다. 이후 미팅에서 모순 그룹의 모든 행을 제외하고 재검증하기로 결정했습니다. CIC2018 EXP47 완료: 실제 expert 호출·채택 0건. ToN EXP48도 같은 정제 기준으로 실행합니다. 상세 진단·시각화·최신 실행 상태는 통합본 목차의 「모순 제거 후 구조 진단 · CIC2018와 ToN」(09-16) 을 확인하세요.

imbalcic · 데이터 감사 · 2026-09-11 (09-15 개정) · 코드 docs/\allowbreak{}research/\allowbreak{}20260911/\allowbreak{}dataset\_\allowbreak{}quality\_\allowbreak{}audit/\allowbreak{}

09-22 갱신: §2에서 CIC·ToN의 다섯 조건, 정답 없는 고정 영역 배정, 점수 차이의 원인, 지표별 검증 목적과 추가 실험를 제시한다(EXP56). 아래 데이터 감사와 과거 참조값은 원래 평가 집합을 명시해 보존한다. 전 행(66,935,021 × 46)을 hash로 묶어 (1) 모순 (한 벡터가 두 라벨)· 중복 (같은 벡터·라벨의 반복)· 누수 (test 벡터가 train에 존재)가 클래스별로 어떻게 생겼는지, (2) 모순 제약과 bank 정답 선택 진단이 각각 무엇인지, (3) 다른 연구는 이 셋을 어떻게 처리했는지, (4) 중복과 모순을 분리해 처리하면 방법별·클래스별로 무엇이 달라지는지, (5) 시간순으로 어떤 클래스가 언제 나타나는지를 정리한다.

\subsection*{1. 문제: 모순 · 중복 · 누수}
세 가지를 구분해서 부른다. 지금까지 "동일 벡터·상반 라벨"과 "중복"을 섞어 쓴 탓에 처리 실험의 조건이 흐려졌다.

\begingroup
\footnotesize
\begin{longtable}{>{\raggedright\arraybackslash}p{\dimexpr 0.21396\linewidth-2\tabcolsep\relax}>{\raggedright\arraybackslash}p{\dimexpr 0.38288\linewidth-2\tabcolsep\relax}>{\raggedright\arraybackslash}p{\dimexpr 0.40315\linewidth-2\tabcolsep\relax}}
\toprule
이름 & 정의 & 성질 \\ \midrule
\endfirsthead
이름 & 정의 & 성질 \\ \midrule
\endhead
{}모순 & 한 feature 벡터가 서로 다른 라벨을 가짐 & 데이터셋 고유. 어떤 분할을 써도 남는다 \\
{}중복 & 같은 (feature 벡터, 라벨) 쌍이 2회 이상 나타남 & 데이터셋 고유. 공격이 반복되면 정상적으로 생긴다 \\
{}누수 & test의 벡터가 train에도 존재함 & 분할에 따라 달라진다. 데이터의 성질이 아니라 평가 설계의 결과 \\
\bottomrule
\end{longtable}
\endgroup
두 정의는 직교한다. 모순 벡터도 중복될 수 있고, 중복이 전혀 없는 모순 벡터도 있다. 누수는 층위가 다르다 — 같은 데이터라도 랜덤 split이면 크고 시간순 split이면 작다. 아래 §4의 처리 실험은 이 직교성 위에서 설계해야 한다. 중복은 따로 도표화하지 않는다 — volumetric 공격이 정적 표적을 때리면 거의 같은 flow가 수백만 개 나오는 것이 정상이고(Flood et al. EuroS\&P 2024), 그 반복은 공격의 서명이자 사전확률이다. 지워야 할 오류가 아니므로 §4에서 처리 축으로만 쓴다.

\subsection*{1a. 모순 — 어느 클래스끼리 벡터를 나눠 갖는가}
NetFlow 한 행은 46개 feature로 요약된 흐름이다. IDS 데이터셋의 라벨은 패킷 내용이 아니라 공격자 IP·대상 포트·시간 창 으로 붙는다(Liu et al. CNS'22; Engelen et al. WTMC'21). 그래서 창 안의 평범한 흐름이 "공격"이 되고, 창 밖의 같은 흐름이 "benign"이 된다. A를 B로 잘못 적었다는 뜻이 아니라, 같은 벡터가 두 라벨을 동시에 갖는다는 뜻이다.

\subsection*{CIC-IDS2018 NF-CSE-CIC-IDS2018-v3 20,115,529행 · 고유 벡터 8,805,692 · 모순 그룹 20,609개 · 모순 벡터에 속한 행 26.6\%}
\begingroup
\footnotesize
\begin{longtable}{>{\raggedright\arraybackslash}p{\dimexpr 0.36551\linewidth-2\tabcolsep\relax}>{\raggedright\arraybackslash}p{\dimexpr 0.23892\linewidth-2\tabcolsep\relax}>{\raggedright\arraybackslash}p{\dimexpr 0.22627\linewidth-2\tabcolsep\relax}>{\raggedright\arraybackslash}p{\dimexpr 0.16930\linewidth-2\tabcolsep\relax}}
\toprule
같은 벡터에 붙은 라벨 A /\allowbreak{} B & A 행 수 (A 전체 대비) & B 행 수 (B 전체 대비) & 벡터 수 \\ \midrule
\endfirsthead
같은 벡터에 붙은 라벨 A /\allowbreak{} B & A 행 수 (A 전체 대비) & B 행 수 (B 전체 대비) & 벡터 수 \\ \midrule
\endhead
{}benign /\allowbreak{} infiltration 시나리오: benign 4,676,451 /\allowbreak{} infilteration 141,386 기간: 02-14–03-02 /\allowbreak{} 02-28–03-01 & 4,676,451 (26.7\%) & 141,386 (75.1\%) & 20,378 \\
{}brute\_\allowbreak{}force /\allowbreak{} dos 시나리오: ftp\_\allowbreak{}bruteforce 386,720 · ssh\_\allowbreak{}bruteforce 60 /\allowbreak{} dos\_\allowbreak{}attacks\_\allowbreak{}slowhttptest 105,550 기간: 02-14–02-14 /\allowbreak{} 02-16–02-16 & 386,780 (67.2\%) & 105,550 (34.8\%) & 1 \\
{}benign /\allowbreak{} web\_\allowbreak{}attacks 시나리오: benign 44,528 /\allowbreak{} brute\_\allowbreak{}force\_\allowbreak{}web 984 · brute\_\allowbreak{}force\_\allowbreak{}xss 436 · sql\_\allowbreak{}injection 344 기간: 02-14–03-02 /\allowbreak{} 02-22–02-22 & 44,528 (0.3\%) & 1,764 (69.5\%) & 298 \\
{}benign /\allowbreak{} ddos 시나리오: benign 936 /\allowbreak{} ddos\_\allowbreak{}attacks\_\allowbreak{}loic\_\allowbreak{}http 8,325 기간: 02-14–03-02 /\allowbreak{} 02-20–02-20 & 936 (0.0\%) & 8,325 (0.6\%) & 1 \\
\bottomrule
\end{longtable}
\endgroup
\subsection*{ToN-IoT NF-ToN-IoT-v3 27,520,260행 · 고유 벡터 4,087,189 · 모순 그룹 41,422개 · 모순 벡터에 속한 행 60.7\%}
\begingroup
\footnotesize
\begin{longtable}{>{\raggedright\arraybackslash}p{\dimexpr 0.35957\linewidth-2\tabcolsep\relax}>{\raggedright\arraybackslash}p{\dimexpr 0.23611\linewidth-2\tabcolsep\relax}>{\raggedright\arraybackslash}p{\dimexpr 0.23611\linewidth-2\tabcolsep\relax}>{\raggedright\arraybackslash}p{\dimexpr 0.16821\linewidth-2\tabcolsep\relax}}
\toprule
같은 벡터에 붙은 라벨 A /\allowbreak{} B & A 행 수 (A 전체 대비) & B 행 수 (B 전체 대비) & 벡터 수 \\ \midrule
\endfirsthead
같은 벡터에 붙은 라벨 A /\allowbreak{} B & A 행 수 (A 전체 대비) & B 행 수 (B 전체 대비) & 벡터 수 \\ \midrule
\endhead
{}benign /\allowbreak{} scanning 시나리오: benign 8,380,313 /\allowbreak{} scanning 1,319,888 기간: 04-23–04-29 /\allowbreak{} 04-23–04-24 & 8,380,313 (49.9\%) & 1,319,888 (97.1\%) & 8,454 \\
{}benign /\allowbreak{} ddos 시나리오: benign 2,305,922 /\allowbreak{} ddos 3,160,750 기간: 04-23–04-29 /\allowbreak{} 04-25–04-26 & 2,305,922 (13.7\%) & 3,160,750 (76.3\%) & 1,682 \\
{}ddos /\allowbreak{} xss 시나리오: ddos 2,862,347 /\allowbreak{} xss 344,748 기간: 04-25–04-26 /\allowbreak{} 04-27–04-27 & 2,862,347 (69.1\%) & 344,748 (12.2\%) & 1,888 \\
{}ddos /\allowbreak{} injection 시나리오: ddos 3,107,486 /\allowbreak{} injection 13,028 기간: 04-25–04-26 /\allowbreak{} 04-25–04-25 & 3,107,486 (75.0\%) & 13,028 (3.4\%) & 727 \\
\bottomrule
\end{longtable}
\endgroup
\subsection*{BoT-IoT NF-BoT-IoT-v3 16,933,808행 · 고유 벡터 669,506 · 모순 그룹 722개 · 모순 벡터에 속한 행 6.1\%}
\begingroup
\footnotesize
\begin{longtable}{>{\raggedright\arraybackslash}p{\dimexpr 0.37500\linewidth-2\tabcolsep\relax}>{\raggedright\arraybackslash}p{\dimexpr 0.22368\linewidth-2\tabcolsep\relax}>{\raggedright\arraybackslash}p{\dimexpr 0.23026\linewidth-2\tabcolsep\relax}>{\raggedright\arraybackslash}p{\dimexpr 0.17105\linewidth-2\tabcolsep\relax}}
\toprule
같은 벡터에 붙은 라벨 A /\allowbreak{} B & A 행 수 (A 전체 대비) & B 행 수 (B 전체 대비) & 벡터 수 \\ \midrule
\endfirsthead
같은 벡터에 붙은 라벨 A /\allowbreak{} B & A 행 수 (A 전체 대비) & B 행 수 (B 전체 대비) & 벡터 수 \\ \midrule
\endhead
{}benign /\allowbreak{} reconnaissance 시나리오: benign 33,189 /\allowbreak{} reconnaissance 603,218 기간: 05-15–05-22 /\allowbreak{} 05-15–05-22 & 33,189 (63.8\%) & 603,218 (35.6\%) & 463 \\
{}ddos /\allowbreak{} dos 시나리오: ddos 114,050 /\allowbreak{} dos 158,927 기간: 06-04–06-04 /\allowbreak{} 06-04–06-04 & 114,050 (1.6\%) & 158,927 (2.0\%) & 252 \\
{}dos /\allowbreak{} reconnaissance 시나리오: dos 99,697 /\allowbreak{} reconnaissance 42,844 기간: 06-04–06-04 /\allowbreak{} 05-15–05-22 & 99,697 (1.2\%) & 42,844 (2.5\%) & 16 \\
{}benign /\allowbreak{} dos 시나리오: benign 4,005 /\allowbreak{} dos 99,687 기간: 05-15–05-22 /\allowbreak{} 06-04–06-04 & 4,005 (7.7\%) & 99,687 (1.2\%) & 14 \\
\bottomrule
\end{longtable}
\endgroup
\subsection*{UNSW-NB15 NF-UNSW-NB15-v3 2,365,424행 · 고유 벡터 1,061,597 · 모순 그룹 1,050개 · 모순 벡터에 속한 행 1.3\%}
\begingroup
\footnotesize
\begin{longtable}{>{\raggedright\arraybackslash}p{\dimexpr 0.37833\linewidth-2\tabcolsep\relax}>{\raggedright\arraybackslash}p{\dimexpr 0.22500\linewidth-2\tabcolsep\relax}>{\raggedright\arraybackslash}p{\dimexpr 0.22500\linewidth-2\tabcolsep\relax}>{\raggedright\arraybackslash}p{\dimexpr 0.17167\linewidth-2\tabcolsep\relax}}
\toprule
같은 벡터에 붙은 라벨 A /\allowbreak{} B & A 행 수 (A 전체 대비) & B 행 수 (B 전체 대비) & 벡터 수 \\ \midrule
\endfirsthead
같은 벡터에 붙은 라벨 A /\allowbreak{} B & A 행 수 (A 전체 대비) & B 행 수 (B 전체 대비) & 벡터 수 \\ \midrule
\endhead
{}backdoor /\allowbreak{} generic 시나리오: backdoor 891 /\allowbreak{} generic 13,092 기간: 02-18–02-18 /\allowbreak{} 01-22–02-18 & 891 (19.1\%) & 13,092 (66.6\%) & 43 \\
{}fuzzers /\allowbreak{} reconnaissance 시나리오: fuzzers 5,333 /\allowbreak{} reconnaissance 5,300 기간: 01-22–02-18 /\allowbreak{} 01-22–02-18 & 5,333 (15.8\%) & 5,300 (31.0\%) & 251 \\
{}exploits /\allowbreak{} fuzzers 시나리오: exploits 2,550 /\allowbreak{} fuzzers 6,079 기간: 01-22–02-18 /\allowbreak{} 01-22–02-18 & 2,550 (6.0\%) & 6,079 (18.0\%) & 479 \\
{}exploits /\allowbreak{} reconnaissance 시나리오: exploits 2,026 /\allowbreak{} reconnaissance 5,215 기간: 01-22–02-18 /\allowbreak{} 01-22–02-18 & 2,026 (4.7\%) & 5,215 (30.5\%) & 158 \\
\bottomrule
\end{longtable}
\endgroup
\subsection*{1b. 누수 — test 행을 암기·모순·신규로 나누면}
파이프라인과 같은 시간순 분할에서 test 행을 셋으로 나눈다: train에 같은 라벨 로 있던 벡터(암기하면 맞는다), train에 다른 라벨 로 있던 벡터(모순이라 무엇을 답해도 한쪽은 틀린다), 처음 보는 벡터(여기서만 일반화가 측정된다).

같은 데이터를 랜덤 80/\allowbreak{}20으로 자르면 test 벡터의 59–99\%가 train에 그대로 남아, 모델 없이 조회표만 써도 정확도가 84–98\%까지 나온다. 누수는 데이터의 성질이 아니라 분할이 만든다. 위 그림은 파이프라인과 같은 시간순 분할 기준이다.

\subsection*{2. 다섯 조건 비교와 expert 생성 품질 진단}
\subsection*{2.1. 비교 기준과 평가 설계}
Expert가 입력을 보고 클래스를 구분하는지 평가하려면, 정답은 채점에만 사용하고 expert 선택이나 오답 복구에 제공하지 않아야 한다. 아래 다섯 조건은 같은 bank와 test를 사용하며, 이전 2번의 정답 클래스 배정을 정답을 쓰지 않는 고정 영역 배정으로 교체했다.

\begingroup
\footnotesize
\begin{longtable}{>{\raggedright\arraybackslash}p{\dimexpr 0.27778\linewidth-2\tabcolsep\relax}>{\raggedright\arraybackslash}p{\dimexpr 0.37778\linewidth-2\tabcolsep\relax}>{\raggedright\arraybackslash}p{\dimexpr 0.34444\linewidth-2\tabcolsep\relax}}
\toprule
조건 & 예측 규칙 & 무엇을 확인하는가 \\ \midrule
\endfirsthead
조건 & 예측 규칙 & 무엇을 확인하는가 \\ \midrule
\endhead
{}1. Global & 모든 행에 global의 원래 예측 & 공통 기준선 \\
{}2. 고정 영역 배정 & 특징 표현·global 확률의 최근접 centroid로 expert를 배정하고 원래 답을 사용 & 정해진 입력 영역과 expert의 결합 성능; S/\allowbreak{}V 없음 \\
{}3. 각 expert 전체 test & 각 expert를 전체 test의 모든 클래스에 독립 적용 & 해당 클래스를 맞히는 능력과 다른 클래스를 오인하는 정도 \\
{}4. Scorer-only & 기존 S+\allowbreak{}V의 scorer 후보·호출 조건 유지, verifier만 제거 & 동일 호출 조건의 scorer-only 결과 \\
{}5. Scorer +\allowbreak{} verifier & 기존 scorer·호출 조건·verifier 사용 & 제안 모델의 최종 성능 \\
\bottomrule
\end{longtable}
\endgroup
진단은 세 질문으로 나눈다. ① Global보다 오탐을 억제하면서 더 잘 구분하는가? ② 사전에 정한 입력 영역에서 global·다른 expert·상수 예측보다 나은가? ③ 그 차이가 context 생성 방식의 효과이며 독립 평가에서도 재현되는가? 각 자료가 답하는 질문과 추가 대조를 §4–6에 연결한다.

생성 조건: 1 \ensuremath{\leq} K < C. K는 global을 제외한 expert 수, C는 benign을 포함한 모델 클래스 수다. 생성 코드에서 K\ensuremath{\geq}C 후보를 제외하도록 적용했다. 정답 클래스 배정의 문제는 K를 줄이는 것만으로 해결되지 않으므로 그 배정 규칙도 주 평가에서 제거했다. 고정 영역 배정은 관측 가능한 규칙의 baseline이며 oracle이나 성능 상한으로 부르지 않는다.

현재 수치는 기존 bank를 이용한 재진단이다. CIC K=\allowbreak{}8·C=\allowbreak{}7은 새 제약을 충족하지 않으므로 새 설계의 결과로 주장하지 않는다. ToN K=\allowbreak{}7·C=\allowbreak{}10은 개수 조건을 충족한다. 이번에 bank 재학습은 하지 않았다.

\subsection*{2.2. 클래스별 점수 차이}
\subsection*{CIC infiltration: 이전 정답 배정에서만 F1 0.9140이 나온 이유}
\begingroup
\footnotesize
\begin{longtable}{>{\raggedright\arraybackslash}p{\dimexpr 0.18546\linewidth-2\tabcolsep\relax}>{\raggedright\arraybackslash}p{\dimexpr 0.13283\linewidth-2\tabcolsep\relax}>{\raggedright\arraybackslash}p{\dimexpr 0.13283\linewidth-2\tabcolsep\relax}>{\raggedright\arraybackslash}p{\dimexpr 0.13283\linewidth-2\tabcolsep\relax}>{\raggedright\arraybackslash}p{\dimexpr 0.15038\linewidth-2\tabcolsep\relax}>{\raggedright\arraybackslash}p{\dimexpr 0.13283\linewidth-2\tabcolsep\relax}>{\raggedright\arraybackslash}p{\dimexpr 0.13283\linewidth-2\tabcolsep\relax}}
\toprule
예측 규칙 & TP & FP & FN & Precision & Recall & F1 \\ \midrule
\endfirsthead
예측 규칙 & TP & FP & FN & Precision & Recall & F1 \\ \midrule
\endhead
{}Global & 2,733 & 3,948 & 6,009 & 0.4091 & 0.3126 & 0.3544 \\
{}e6 전체 test & 7,556 & 806,387 & 1,186 & 0.0093 & 0.8643 & 0.0184 \\
{}제외한 정답 클래스 배정 & 7,556 & 236 & 1,186 & 0.9697 & 0.8643 & 0.9140 \\
\bottomrule
\end{longtable}
\endgroup
정답 배정은 실제 infiltration 행만 e6에 주고 다른 클래스는 다른 예측기에 보낸다. 그래서 TP 7,556과 FN 1,186은 e6 그대로인데, e6 전체 적용에서 생기던 FP 806,387 대신 선택된 다른 예측기의 FP 236만 남는다. F1 =\allowbreak{} 2×7,556 /\allowbreak{} (2×7,556 +\allowbreak{} 236 +\allowbreak{} 1,186) =\allowbreak{} 0.9140이다. 정답을 아는 배정 규칙이 음성을 걸러준 결과이며 e6의 분류 성능이 0.9140이라는 뜻이 아니다. 이 이유로 해당 값을 주 비교에서 제거했다.

\subsection*{ToN scanning: 개별 expert보다 S+\allowbreak{}V의 F1이 높은 이유}
\begingroup
\footnotesize
\begin{longtable}{>{\raggedright\arraybackslash}p{\dimexpr 0.26157\linewidth-2\tabcolsep\relax}>{\raggedright\arraybackslash}p{\dimexpr 0.12072\linewidth-2\tabcolsep\relax}>{\raggedright\arraybackslash}p{\dimexpr 0.12072\linewidth-2\tabcolsep\relax}>{\raggedright\arraybackslash}p{\dimexpr 0.12072\linewidth-2\tabcolsep\relax}>{\raggedright\arraybackslash}p{\dimexpr 0.13481\linewidth-2\tabcolsep\relax}>{\raggedright\arraybackslash}p{\dimexpr 0.12072\linewidth-2\tabcolsep\relax}>{\raggedright\arraybackslash}p{\dimexpr 0.12072\linewidth-2\tabcolsep\relax}}
\toprule
예측 규칙 & TP & FP & FN & Precision & Recall & F1 \\ \midrule
\endfirsthead
예측 규칙 & TP & FP & FN & Precision & Recall & F1 \\ \midrule
\endhead
{}Global & 452 & 50 & 5,577 & 0.9004 & 0.0750 & 0.1384 \\
{}e2 전체 test & 5,439 & 54,321 & 590 & 0.0910 & 0.9021 & 0.1653 \\
{}e4 전체 test: 개별 expert 최고 F1 & 4,240 & 7,715 & 1,789 & 0.3547 & 0.7033 & 0.4715 \\
{}S-only & 4,270 & 287 & 1,759 & 0.9370 & 0.7082 & 0.8067 \\
{}S+\allowbreak{}V & 4,245 & 172 & 1,784 & 0.9611 & 0.7041 & 0.8128 \\
\bottomrule
\end{longtable}
\endgroup
S+\allowbreak{}V는 expert 하나를 모든 행에 적용하지 않는다. 입력에 따라 expert를 선택하고 호출·승인 조건을 통과한 답만 사용한다. Scanning TP 4,245개를 확보하면서 FP를 172개로 제한 하므로 F1 =\allowbreak{} 2×4,245 /\allowbreak{} (2×4,245 +\allowbreak{} 172 +\allowbreak{} 1,784) =\allowbreak{} 0.8128이다. 전체 test에 e4를 적용한 경우와 TP는 비슷하지만 FP는 7,715개다. 따라서 합성 정책의 F1은 각 expert 전면 적용 F1의 최댓값보다 높을 수 있다.

실제 예측 출처도 확인했다. S+\allowbreak{}V scanning의 TP는 global 452개 +\allowbreak{} 선택된 expert들 3,793개이며, FP는 global 50개 +\allowbreak{} expert들 122개다. 이 예측을 구성할 때 test 정답은 사용하지 않았고, 정답 배정의 음성 제거와는 평가 규칙이 다르다.

\begingroup
\footnotesize
\begin{longtable}{>{\raggedright\arraybackslash}p{\dimexpr 0.33333\linewidth-2\tabcolsep\relax}>{\raggedright\arraybackslash}p{\dimexpr 0.33333\linewidth-2\tabcolsep\relax}>{\raggedright\arraybackslash}p{\dimexpr 0.33333\linewidth-2\tabcolsep\relax}}
\toprule
예측 출처 & TP & FP \\ \midrule
\endfirsthead
예측 출처 & TP & FP \\ \midrule
\endhead
{}global & 452 & 50 \\
{}expert1 & 436 & 5 \\
{}expert2 & 2,852 & 105 \\
{}expert3 & 8 & 5 \\
{}expert4 & 156 & 6 \\
{}expert5 & 194 & 1 \\
{}expert6 & 71 & 0 \\
{}expert7 & 76 & 0 \\
\bottomrule
\end{longtable}
\endgroup
\subsection*{2.3. 수정된 다섯 조건의 성능}
2번은 정답을 쓰지 않는 고정 영역 배정으로 재계산했다. 1·3·4·5번은 기존 예측·점수와 일치한다. 굵은 값은 다섯 조건의 최고값이며 소수 넷째 자리 표시가 같으면 함께 강조한다. 3번은 K개 독립 실행의 최소–최대이고, 괄호는 최대값 expert다.

\begingroup
\footnotesize
\begin{longtable}{>{\raggedright\arraybackslash}p{\dimexpr 0.14881\linewidth-2\tabcolsep\relax}>{\raggedright\arraybackslash}p{\dimexpr 0.13690\linewidth-2\tabcolsep\relax}>{\raggedright\arraybackslash}p{\dimexpr 0.14881\linewidth-2\tabcolsep\relax}>{\raggedright\arraybackslash}p{\dimexpr 0.19643\linewidth-2\tabcolsep\relax}>{\raggedright\arraybackslash}p{\dimexpr 0.16667\linewidth-2\tabcolsep\relax}>{\raggedright\arraybackslash}p{\dimexpr 0.20238\linewidth-2\tabcolsep\relax}}
\toprule
데이터 & 1. Global & 2. 고정 영역 배정 & 3. 각 expert 전체 test & 4. Scorer-only & 5. Scorer +\allowbreak{} verifier \\ \midrule
\endfirsthead
데이터 & 1. Global & 2. 고정 영역 배정 & 3. 각 expert 전체 test & 4. Scorer-only & 5. Scorer +\allowbreak{} verifier \\ \midrule
\endhead
{}CIC-IDS2018 & 0.8107 & 0.7477 & 0.7001– 0.7946 (e2) & 0.8107 & 0.8107 \\
{}ToN-IoT & 0.6837 & 0.6620 & 0.6205– 0.6859 (e7) & 0.7549 & 0.7555 \\
\bottomrule
\end{longtable}
\endgroup
\subsection*{CIC-IDS2018 — 클래스별 F1}
\begingroup
\footnotesize
\begin{longtable}{>{\raggedright\arraybackslash}p{\dimexpr 0.15362\linewidth-2\tabcolsep\relax}>{\raggedright\arraybackslash}p{\dimexpr 0.13307\linewidth-2\tabcolsep\relax}>{\raggedright\arraybackslash}p{\dimexpr 0.13307\linewidth-2\tabcolsep\relax}>{\raggedright\arraybackslash}p{\dimexpr 0.12622\linewidth-2\tabcolsep\relax}>{\raggedright\arraybackslash}p{\dimexpr 0.20157\linewidth-2\tabcolsep\relax}>{\raggedright\arraybackslash}p{\dimexpr 0.13307\linewidth-2\tabcolsep\relax}>{\raggedright\arraybackslash}p{\dimexpr 0.11937\linewidth-2\tabcolsep\relax}}
\toprule
클래스 & Test 수 & 1. Global & 2. 고정 영역 & 3. Expert 범위 & 4. S-only & 5. S+\allowbreak{}V \\ \midrule
\endfirsthead
클래스 & Test 수 & 1. Global & 2. 고정 영역 & 3. Expert 범위 & 4. S-only & 5. S+\allowbreak{}V \\ \midrule
\endhead
{}benign & 2,652,203 & 0.9978 & 0.9838 & 0.8199– 0.9975 (e2) & 0.9978 & 0.9978 \\
{}ddos & 262,680 & 0.9991 & 0.9937 & 0.9919– 1.0000 (e3) & 0.9991 & 0.9991 \\
{}bot & 41,542 & 0.9982 & 0.9977 & 0.9934– 0.9978 (e2) & 0.9982 & 0.9982 \\
{}dos & 39,483 & 0.9920 & 0.9920 & 0.9846– 0.9923 (e7) & 0.9920 & 0.9920 \\
{}brute\_\allowbreak{}force & 37,696 & 0.9997 & 0.9997 & 0.9997– 0.9997 (e1) & 0.9997 & 0.9997 \\
{}infiltration & 8,742 & 0.3544 & 0.0862 & 0.0184– 0.2209 (e2) & 0.3544 & 0.3544 \\
{}web\_\allowbreak{}attacks & 127 & 0.3333 & 0.1808 & 0.0152– 0.3622 (e2) & 0.3333 & 0.3333 \\
\bottomrule
\end{longtable}
\endgroup
\begin{landscape}
\begingroup
\scriptsize
\begin{longtable}{>{\raggedright\arraybackslash}p{\dimexpr 0.09232\linewidth-2\tabcolsep\relax}>{\raggedright\arraybackslash}p{\dimexpr 0.07817\linewidth-2\tabcolsep\relax}>{\raggedright\arraybackslash}p{\dimexpr 0.06873\linewidth-2\tabcolsep\relax}>{\raggedright\arraybackslash}p{\dimexpr 0.07345\linewidth-2\tabcolsep\relax}>{\raggedright\arraybackslash}p{\dimexpr 0.06873\linewidth-2\tabcolsep\relax}>{\raggedright\arraybackslash}p{\dimexpr 0.06873\linewidth-2\tabcolsep\relax}>{\raggedright\arraybackslash}p{\dimexpr 0.06873\linewidth-2\tabcolsep\relax}>{\raggedright\arraybackslash}p{\dimexpr 0.06873\linewidth-2\tabcolsep\relax}>{\raggedright\arraybackslash}p{\dimexpr 0.06873\linewidth-2\tabcolsep\relax}>{\raggedright\arraybackslash}p{\dimexpr 0.06873\linewidth-2\tabcolsep\relax}>{\raggedright\arraybackslash}p{\dimexpr 0.06873\linewidth-2\tabcolsep\relax}>{\raggedright\arraybackslash}p{\dimexpr 0.06873\linewidth-2\tabcolsep\relax}>{\raggedright\arraybackslash}p{\dimexpr 0.06873\linewidth-2\tabcolsep\relax}>{\raggedright\arraybackslash}p{\dimexpr 0.06873\linewidth-2\tabcolsep\relax}}
\toprule
클래스 & Test 수 & 1. G & 2. 고정 영역 & 3. e1 & 3. e2 & 3. e3 & 3. e4 & 3. e5 & 3. e6 & 3. e7 & 3. e8 & 4. S & 5. S+\allowbreak{}V \\ \midrule
\endfirsthead
클래스 & Test 수 & 1. G & 2. 고정 영역 & 3. e1 & 3. e2 & 3. e3 & 3. e4 & 3. e5 & 3. e6 & 3. e7 & 3. e8 & 4. S & 5. S+\allowbreak{}V \\ \midrule
\endhead
{}benign & 2,652,203 & 0.9978 & 0.9838 & 0.8948 & 0.9975 & 0.9042 & 0.9027 & 0.8954 & 0.8199 & 0.8925 & 0.8962 & 0.9978 & 0.9978 \\
{}ddos & 262,680 & 0.9991 & 0.9937 & 0.9996 & 0.9919 & 1.0000 & 1.0000 & 0.9983 & 1.0000 & 0.9998 & 1.0000 & 0.9991 & 0.9991 \\
{}bot & 41,542 & 0.9982 & 0.9977 & 0.9934 & 0.9978 & 0.9978 & 0.9974 & 0.9973 & 0.9974 & 0.9973 & 0.9973 & 0.9982 & 0.9982 \\
{}dos & 39,483 & 0.9920 & 0.9920 & 0.9846 & 0.9921 & 0.9919 & 0.9913 & 0.9887 & 0.9903 & 0.9923 & 0.9920 & 0.9920 & 0.9920 \\
{}brute\_\allowbreak{}force & 37,696 & 0.9997 & 0.9997 & 0.9997 & 0.9997 & 0.9997 & 0.9997 & 0.9997 & 0.9997 & 0.9997 & 0.9997 & 0.9997 & 0.9997 \\
{}infiltration & 8,742 & 0.3544 & 0.0862 & 0.0270 & 0.2209 & 0.0262 & 0.0261 & 0.0271 & 0.0184 & 0.0235 & 0.0256 & 0.3544 & 0.3544 \\
{}web\_\allowbreak{}attacks & 127 & 0.3333 & 0.1808 & 0.0152 & 0.3622 & 0.0715 & 0.0695 & 0.0419 & 0.0749 & 0.0536 & 0.1854 & 0.3333 & 0.3333 \\
\bottomrule
\end{longtable}
\endgroup
\end{landscape}
\subsection*{ToN-IoT — 클래스별 F1}
\begingroup
\footnotesize
\begin{longtable}{>{\raggedright\arraybackslash}p{\dimexpr 0.14185\linewidth-2\tabcolsep\relax}>{\raggedright\arraybackslash}p{\dimexpr 0.13481\linewidth-2\tabcolsep\relax}>{\raggedright\arraybackslash}p{\dimexpr 0.13481\linewidth-2\tabcolsep\relax}>{\raggedright\arraybackslash}p{\dimexpr 0.12777\linewidth-2\tabcolsep\relax}>{\raggedright\arraybackslash}p{\dimexpr 0.20523\linewidth-2\tabcolsep\relax}>{\raggedright\arraybackslash}p{\dimexpr 0.13481\linewidth-2\tabcolsep\relax}>{\raggedright\arraybackslash}p{\dimexpr 0.12072\linewidth-2\tabcolsep\relax}}
\toprule
클래스 & Test 수 & 1. Global & 2. 고정 영역 & 3. Expert 범위 & 4. S-only & 5. S+\allowbreak{}V \\ \midrule
\endfirsthead
클래스 & Test 수 & 1. Global & 2. 고정 영역 & 3. Expert 범위 & 4. S-only & 5. S+\allowbreak{}V \\ \midrule
\endhead
{}benign & 1,474,302 & 0.9119 & 0.8913 & 0.8250– 0.8932 (e6) & 0.9199 & 0.9200 \\
{}xss & 343,854 & 0.8070 & 0.8574 & 0.7309– 0.8616 (e5) & 0.8372 & 0.8374 \\
{}password & 210,951 & 0.8909 & 0.8921 & 0.8835– 0.8922 (e2) & 0.8909 & 0.8909 \\
{}ddos & 103,338 & 0.9687 & 0.9653 & 0.9626– 0.9746 (e4) & 0.9727 & 0.9727 \\
{}injection & 68,891 & 0.9315 & 0.9027 & 0.3883– 0.9539 (e7) & 0.9321 & 0.9321 \\
{}backdoor & 40,030 & 0.9987 & 0.9985 & 0.9973– 0.9989 (e7) & 0.9988 & 0.9988 \\
{}dos & 22,458 & 0.9831 & 0.7312 & 0.4594– 0.9766 (e7) & 0.9839 & 0.9832 \\
{}scanning & 6,029 & 0.1384 & 0.3062 & 0.1426– 0.4715 (e4) & 0.8067 & 0.8128 \\
{}mitm & 1,161 & 0.0674 & 0.0432 & 0.0282– 0.0614 (e3) & 0.0673 & 0.0674 \\
{}ransomware & 709 & 0.1397 & 0.0319 & 0.0305– 0.1237 (e2) & 0.1396 & 0.1397 \\
\bottomrule
\end{longtable}
\endgroup
\begin{landscape}
\begingroup
\scriptsize
\begin{longtable}{>{\raggedright\arraybackslash}p{\dimexpr 0.09001\linewidth-2\tabcolsep\relax}>{\raggedright\arraybackslash}p{\dimexpr 0.08485\linewidth-2\tabcolsep\relax}>{\raggedright\arraybackslash}p{\dimexpr 0.07454\linewidth-2\tabcolsep\relax}>{\raggedright\arraybackslash}p{\dimexpr 0.07970\linewidth-2\tabcolsep\relax}>{\raggedright\arraybackslash}p{\dimexpr 0.07454\linewidth-2\tabcolsep\relax}>{\raggedright\arraybackslash}p{\dimexpr 0.07454\linewidth-2\tabcolsep\relax}>{\raggedright\arraybackslash}p{\dimexpr 0.07454\linewidth-2\tabcolsep\relax}>{\raggedright\arraybackslash}p{\dimexpr 0.07454\linewidth-2\tabcolsep\relax}>{\raggedright\arraybackslash}p{\dimexpr 0.07454\linewidth-2\tabcolsep\relax}>{\raggedright\arraybackslash}p{\dimexpr 0.07454\linewidth-2\tabcolsep\relax}>{\raggedright\arraybackslash}p{\dimexpr 0.07454\linewidth-2\tabcolsep\relax}>{\raggedright\arraybackslash}p{\dimexpr 0.07454\linewidth-2\tabcolsep\relax}>{\raggedright\arraybackslash}p{\dimexpr 0.07454\linewidth-2\tabcolsep\relax}}
\toprule
클래스 & Test 수 & 1. G & 2. 고정 영역 & 3. e1 & 3. e2 & 3. e3 & 3. e4 & 3. e5 & 3. e6 & 3. e7 & 4. S & 5. S+\allowbreak{}V \\ \midrule
\endfirsthead
클래스 & Test 수 & 1. G & 2. 고정 영역 & 3. e1 & 3. e2 & 3. e3 & 3. e4 & 3. e5 & 3. e6 & 3. e7 & 4. S & 5. S+\allowbreak{}V \\ \midrule
\endhead
{}benign & 1,474,302 & 0.9119 & 0.8913 & 0.8878 & 0.8856 & 0.8912 & 0.8250 & 0.8901 & 0.8932 & 0.8344 & 0.9199 & 0.9200 \\
{}xss & 343,854 & 0.8070 & 0.8574 & 0.8603 & 0.8448 & 0.8358 & 0.8175 & 0.8616 & 0.8528 & 0.7309 & 0.8372 & 0.8374 \\
{}password & 210,951 & 0.8909 & 0.8921 & 0.8835 & 0.8922 & 0.8920 & 0.8835 & 0.8921 & 0.8921 & 0.8897 & 0.8909 & 0.8909 \\
{}ddos & 103,338 & 0.9687 & 0.9653 & 0.9626 & 0.9697 & 0.9648 & 0.9746 & 0.9656 & 0.9689 & 0.9644 & 0.9727 & 0.9727 \\
{}injection & 68,891 & 0.9315 & 0.9027 & 0.8896 & 0.8520 & 0.8467 & 0.3883 & 0.8473 & 0.9140 & 0.9539 & 0.9321 & 0.9321 \\
{}backdoor & 40,030 & 0.9987 & 0.9985 & 0.9973 & 0.9979 & 0.9988 & 0.9974 & 0.9984 & 0.9983 & 0.9989 & 0.9988 & 0.9988 \\
{}dos & 22,458 & 0.9831 & 0.7312 & 0.9496 & 0.9709 & 0.4594 & 0.7779 & 0.7576 & 0.8835 & 0.9766 & 0.9839 & 0.9832 \\
{}scanning & 6,029 & 0.1384 & 0.3062 & 0.1430 & 0.1653 & 0.1719 & 0.4715 & 0.4447 & 0.1426 & 0.4467 & 0.8067 & 0.8128 \\
{}mitm & 1,161 & 0.0674 & 0.0432 & 0.0319 & 0.0282 & 0.0614 & 0.0385 & 0.0383 & 0.0345 & 0.0331 & 0.0673 & 0.0674 \\
{}ransomware & 709 & 0.1397 & 0.0319 & 0.0906 & 0.1237 & 0.0965 & 0.0310 & 0.0305 & 0.1232 & 0.0308 & 0.1396 & 0.1397 \\
\bottomrule
\end{longtable}
\endgroup
\end{landscape}
\subsection*{2.4. 어떤 자료가 expert의 어떤 역량을 검증하는가?}
핵심은 그 클래스 이름을 많이 출력하는가 가 아니라 그 클래스인 샘플과 아닌 샘플을 구분하는가 다. 아래 자료는 동일한 성능을 중복 나열한 것이 아니라 서로 다른 실패 원인을 구분하는 검사다.

\begingroup
\footnotesize
\begin{longtable}{>{\raggedright\arraybackslash}p{\dimexpr 0.31153\linewidth-2\tabcolsep\relax}>{\raggedright\arraybackslash}p{\dimexpr 0.34424\linewidth-2\tabcolsep\relax}>{\raggedright\arraybackslash}p{\dimexpr 0.34424\linewidth-2\tabcolsep\relax}}
\toprule
자료 & 검증하려는 질문 & 좋은 결과의 의미 /\allowbreak{} 해석 한계 \\ \midrule
\endfirsthead
자료 & 검증하려는 질문 & 좋은 결과의 의미 /\allowbreak{} 해석 한계 \\ \midrule
\endhead
{}전체 test의 precision·recall·F1·FP & 원래 예측 규칙으로 양성과 음성을 구분하는가? & Recall과 precision이 함께 좋아야 한다. 전역 macro-F1만으로 특화 영역의 능력을 판정하지 않는다. \\
{}AP (average precision) & 양성에 음성보다 높은 점수를 부여하는가? & Precision–recall 전반의 순위 성능. 특정 FPR을 고정한 지표는 아니다. 같은 클래스·같은 test에서 G와 비교한다. \\
{}R@FPR\ensuremath{\leq}0.1\% & 음성 1,000개당 최대 1개 오탐 수준에서도 양성을 얼마나 찾는가? & 동일 오탐 한도에서의 식별력 비교. 여기서는 test ROC의 기술적 요약이며 배포 임계값을 정한 결과가 아니다. \\
{}Val에서 고정한 임계값의 test recall·FPR & 다음 데이터에서도 정한 오탐 수준과 탐지율이 유지되는가? & 임계값 전이·안정성 검사. AP가 높아도 test FPR이 크게 벗어나면 그 운영점을 바로 쓸 근거는 없다. \\
{}관측 가능한 영역 내 비교 & 미리 정한 입력 영역에 해당 expert를 두는 것이 적절한가? & 영역과 expert의 적합성 검사. 단일 클래스 영역의 높은 accuracy는 구분 능력을 증명하지 못하므로 상수 대조·영역 macro-F1을 함께 본다. \\
\bottomrule
\end{longtable}
\endgroup
AP와 R@0.1\%는 모두 ①의 식별력 진단에 속하지만 역할이 다르다. AP는 임계값 전반의 precision–recall 관계이고, R@0.1\%가 명시적으로 동일 FPR 한도를 고정한다. 모든 샘플에 같은 점수를 주면 AP는 양성 비율 수준이고, FPR 0.1\%에서는 전부 양성이라고 할 수 없다. 0.1\%는 이번 비교에 사용한 진단 기준이며 실제 허용 오탐률을 확정한 값은 아니다. One-vs-rest 점수로 계산한 AP·R@0.1\%와 원래 다중 클래스 argmax의 F1은 서로 다른 평가다.

\subsection*{CIC-IDS2018 — 양성·음성 구분 능력}
\begin{landscape}
\begingroup
\scriptsize
\begin{longtable}{>{\raggedright\arraybackslash}p{\dimexpr 0.08921\linewidth-2\tabcolsep\relax}>{\raggedright\arraybackslash}p{\dimexpr 0.11325\linewidth-2\tabcolsep\relax}>{\raggedright\arraybackslash}p{\dimexpr 0.08921\linewidth-2\tabcolsep\relax}>{\raggedright\arraybackslash}p{\dimexpr 0.09882\linewidth-2\tabcolsep\relax}>{\raggedright\arraybackslash}p{\dimexpr 0.12286\linewidth-2\tabcolsep\relax}>{\raggedright\arraybackslash}p{\dimexpr 0.09882\linewidth-2\tabcolsep\relax}>{\raggedright\arraybackslash}p{\dimexpr 0.12286\linewidth-2\tabcolsep\relax}>{\raggedright\arraybackslash}p{\dimexpr 0.11806\linewidth-2\tabcolsep\relax}>{\raggedright\arraybackslash}p{\dimexpr 0.14690\linewidth-2\tabcolsep\relax}}
\toprule
Expert & 주로 포함된 클래스 & G AP & Expert AP & 타 expert 최고 AP & G R@0.1\% & Expert R@0.1\% & 타 expert 최고 R & 읽는 법 \\ \midrule
\endfirsthead
Expert & 주로 포함된 클래스 & G AP & Expert AP & 타 expert 최고 AP & G R@0.1\% & Expert R@0.1\% & 타 expert 최고 R & 읽는 법 \\ \midrule
\endhead
{}e1 & bot & 0.9974 & 0.9981 & 0.9996 & 0.9971 & 0.9971 & 0.9971 & 지표 혼재 /\allowbreak{} 이득 제한 \\
{}e2 & benign & 0.9996 & 0.9990 & 0.9996 & 0.6286 & 0.3605 & 0.6613 & G 대비 우위 미확인 \\
{}e3 & brute\_\allowbreak{}force & 1.0000 & 0.9999 & 1.0000 & 1.0000 & 0.9996 & 1.0000 & G 대비 우위 미확인 \\
{}e4 & dos & 0.9909 & 0.9995 & 0.9999 & 0.9849 & 0.9973 & 1.0000 & G 개선 · 타 expert도 유효 \\
{}e5 & brute\_\allowbreak{}force & 1.0000 & 1.0000 & 1.0000 & 1.0000 & 1.0000 & 1.0000 & 지표 혼재 /\allowbreak{} 이득 제한 \\
{}e6 & infiltration & 0.2178 & 0.0944 & 0.1671 & 0.3007 & 0.0991 & 0.1374 & G 대비 우위 미확인 \\
{}e7 & ddos & 1.0000 & 1.0000 & 1.0000 & 1.0000 & 1.0000 & 1.0000 & 지표 혼재 /\allowbreak{} 이득 제한 \\
{}e8 & benign & 0.9996 & 0.9993 & 0.9996 & 0.6286 & 0.4783 & 0.6613 & G 대비 우위 미확인 \\
\bottomrule
\end{longtable}
\endgroup
\end{landscape}
“주로 포함된 클래스”는 context 특화 block의 최다 클래스다. 실제 정답 배정에 사용하지 않는다. 타 expert 최고값은 비교용 요약이며 해당 expert를 test에서 선택한 정책의 점수가 아니다.

특정 클래스 recall이 높을 때 precision과 FP를 함께 확인한다. Precision이 낮고 FP가 많으면 현재 argmax 규칙이 그 클래스를 과잉 출력하는 것이다. 이것만으로 점수 순위까지 나쁘다고 단정하지 않고 위 AP·R@0.1\%와 대조한다.

\begingroup
\footnotesize
\begin{longtable}{>{\raggedright\arraybackslash}p{\dimexpr 0.12446\linewidth-2\tabcolsep\relax}>{\raggedright\arraybackslash}p{\dimexpr 0.16234\linewidth-2\tabcolsep\relax}>{\raggedright\arraybackslash}p{\dimexpr 0.15476\linewidth-2\tabcolsep\relax}>{\raggedright\arraybackslash}p{\dimexpr 0.13961\linewidth-2\tabcolsep\relax}>{\raggedright\arraybackslash}p{\dimexpr 0.12446\linewidth-2\tabcolsep\relax}>{\raggedright\arraybackslash}p{\dimexpr 0.12446\linewidth-2\tabcolsep\relax}>{\raggedright\arraybackslash}p{\dimexpr 0.16991\linewidth-2\tabcolsep\relax}}
\toprule
Expert & 클래스 & 전면 macro-F1 & Precision & Recall & F1 & 다른 클래스에서 온 FP \\ \midrule
\endfirsthead
Expert & 클래스 & 전면 macro-F1 & Precision & Recall & F1 & 다른 클래스에서 온 FP \\ \midrule
\endhead
{}e1 & bot & 0.7020 & 0.9902 & 0.9967 & 0.9934 & 411 \\
{}e2 & benign & 0.7946 & 0.9953 & 0.9998 & 0.9975 & 12,576 \\
{}e3 & brute\_\allowbreak{}force & 0.7130 & 1.0000 & 0.9995 & 0.9997 & 0 \\
{}e4 & dos & 0.7124 & 0.9983 & 0.9845 & 0.9913 & 68 \\
{}e5 & brute\_\allowbreak{}force & 0.7069 & 1.0000 & 0.9995 & 0.9997 & 0 \\
{}e6 & infiltration & 0.7001 & 0.0093 & 0.8643 & 0.0184 & 806,387 \\
{}e7 & ddos & 0.7084 & 1.0000 & 0.9996 & 0.9998 & 5 \\
{}e8 & benign & 0.7280 & 0.9986 & 0.8129 & 0.8962 & 3,118 \\
\bottomrule
\end{longtable}
\endgroup
Val의 음성 점수에서 FPR\ensuremath{\leq}0.1\%를 만족하는 임계값을 고정한 뒤 test에 그대로 적용했다. 실제 test FPR을 먼저 읽고 그때의 recall을 비교한다. 서로 다른 test FPR의 recall만 비교하면 동일 오탐 조건의 우위라고 말할 수 없다.

\begingroup
\footnotesize
\begin{longtable}{>{\raggedright\arraybackslash}p{\dimexpr 0.13000\linewidth-2\tabcolsep\relax}>{\raggedright\arraybackslash}p{\dimexpr 0.16333\linewidth-2\tabcolsep\relax}>{\raggedright\arraybackslash}p{\dimexpr 0.17000\linewidth-2\tabcolsep\relax}>{\raggedright\arraybackslash}p{\dimexpr 0.20333\linewidth-2\tabcolsep\relax}>{\raggedright\arraybackslash}p{\dimexpr 0.15000\linewidth-2\tabcolsep\relax}>{\raggedright\arraybackslash}p{\dimexpr 0.18333\linewidth-2\tabcolsep\relax}}
\toprule
Expert & 클래스 & G test recall & Expert test recall & G test FPR & Expert test FPR \\ \midrule
\endfirsthead
Expert & 클래스 & G test recall & Expert test recall & G test FPR & Expert test FPR \\ \midrule
\endhead
{}e1 & bot & 0.9971 & 0.9971 & 0.177\% & 0.054\% \\
{}e2 & benign & 0.6130 & 0.3968 & 0.090\% & 0.123\% \\
{}e3 & brute\_\allowbreak{}force & 1.0000 & 1.0000 & 0.178\% & 0.762\% \\
{}e4 & dos & 0.9878 & 0.9973 & 0.415\% & 0.091\% \\
{}e5 & brute\_\allowbreak{}force & 1.0000 & 1.0000 & 0.178\% & 0.070\% \\
{}e6 & infiltration & 0.1128 & 0.0746 & 0.003\% & 0.014\% \\
{}e7 & ddos & 1.0000 & 1.0000 & 0.087\% & 0.969\% \\
{}e8 & benign & 0.6130 & 0.5667 & 0.090\% & 0.125\% \\
\bottomrule
\end{longtable}
\endgroup
Test 정답 없이 정한 최근접 centroid 영역에서 비교한다. 상수 대조는 각 영역의 학습용 route 행에서 가장 많은 클래스 하나만 항상 출력한다. 이번 수정에서 추가 계산했다. 영역 내 실제 등장 클래스들의 F1 평균을 표시하며, 같은 행의 세 모델끼리 비교한다.

\begingroup
\footnotesize
\begin{longtable}{>{\raggedright\arraybackslash}p{\dimexpr 0.17793\linewidth-2\tabcolsep\relax}>{\raggedright\arraybackslash}p{\dimexpr 0.14414\linewidth-2\tabcolsep\relax}>{\raggedright\arraybackslash}p{\dimexpr 0.13739\linewidth-2\tabcolsep\relax}>{\raggedright\arraybackslash}p{\dimexpr 0.17117\linewidth-2\tabcolsep\relax}>{\raggedright\arraybackslash}p{\dimexpr 0.20495\linewidth-2\tabcolsep\relax}>{\raggedright\arraybackslash}p{\dimexpr 0.16441\linewidth-2\tabcolsep\relax}}
\toprule
담당 expert /\allowbreak{} 영역 & Test 수 & 등장 클래스 수 & G 영역 macro-F1 & Expert 영역 macro-F1 & 학습 다수 클래스 상수 \\ \midrule
\endfirsthead
담당 expert /\allowbreak{} 영역 & Test 수 & 등장 클래스 수 & G 영역 macro-F1 & Expert 영역 macro-F1 & 학습 다수 클래스 상수 \\ \midrule
\endhead
{}e1 & 41,370 & 2 & 0.5000 & 0.5000 & 0.5000 \\
{}e2 & 2,549,948 & 7 & 0.3157 & 0.1809 & 0.1426 \\
{}e3 & 6,686 & 1 & 1.0000 & 1.0000 & 1.0000 \\
{}e4 & 38,875 & 2 & 0.7999 & 0.6599 & 0.4999 \\
{}e5 & 30,990 & 1 & 1.0000 & 1.0000 & 1.0000 \\
{}e6 & 120,135 & 4 & 0.6181 & 0.4087 & 0.0172 \\
{}e7 & 253,572 & 2 & 0.5000 & 0.9000 & 0.5000 \\
{}e8 & 897 & 5 & 0.1613 & 0.2184 & 0.0462 \\
\bottomrule
\end{longtable}
\endgroup
CIC 영역 3·5는 test에 한 클래스만 있어 상수 예측도 1.0000이다. 이런 영역의 높은 점수는 분류 역량의 증거가 되지 않는다. 영역별 등장 클래스가 달라 이 수치를 서로 다른 영역이나 전체 test macro-F1과 직접 비교하지 않는다. Global·모든 타 expert의 영역 accuracy 원본 행렬도 CSV에 보존했다.

\subsection*{ToN-IoT — 양성·음성 구분 능력}
\begin{landscape}
\begingroup
\scriptsize
\begin{longtable}{>{\raggedright\arraybackslash}p{\dimexpr 0.08987\linewidth-2\tabcolsep\relax}>{\raggedright\arraybackslash}p{\dimexpr 0.10458\linewidth-2\tabcolsep\relax}>{\raggedright\arraybackslash}p{\dimexpr 0.08987\linewidth-2\tabcolsep\relax}>{\raggedright\arraybackslash}p{\dimexpr 0.09967\linewidth-2\tabcolsep\relax}>{\raggedright\arraybackslash}p{\dimexpr 0.12418\linewidth-2\tabcolsep\relax}>{\raggedright\arraybackslash}p{\dimexpr 0.09967\linewidth-2\tabcolsep\relax}>{\raggedright\arraybackslash}p{\dimexpr 0.12418\linewidth-2\tabcolsep\relax}>{\raggedright\arraybackslash}p{\dimexpr 0.11928\linewidth-2\tabcolsep\relax}>{\raggedright\arraybackslash}p{\dimexpr 0.14869\linewidth-2\tabcolsep\relax}}
\toprule
Expert & 주로 포함된 클래스 & G AP & Expert AP & 타 expert 최고 AP & G R@0.1\% & Expert R@0.1\% & 타 expert 최고 R & 읽는 법 \\ \midrule
\endfirsthead
Expert & 주로 포함된 클래스 & G AP & Expert AP & 타 expert 최고 AP & G R@0.1\% & Expert R@0.1\% & 타 expert 최고 R & 읽는 법 \\ \midrule
\endhead
{}e1 & injection & 0.9736 & 0.9604 & 0.9800 & 0.9380 & 0.8023 & 0.9434 & G 대비 우위 미확인 \\
{}e2 & scanning & 0.3310 & 0.8312 & 0.6794 & 0.1415 & 0.7462 & 0.6004 & 두 지표의 관측 우위 \\
{}e3 & dos & 0.9917 & 0.9918 & 0.9921 & 0.9887 & 0.9876 & 0.9863 & 지표 혼재 /\allowbreak{} 이득 제한 \\
{}e4 & injection & 0.9736 & 0.9666 & 0.9800 & 0.9380 & 0.7739 & 0.9434 & G 대비 우위 미확인 \\
{}e5 & ransomware & 0.0477 & 0.2267 & 0.0693 & 0.0310 & 0.6996 & 0.0085 & 두 지표의 관측 우위 \\
{}e6 & benign & 0.9699 & 0.9709 & 0.9749 & 0.3869 & 0.6805 & 0.7753 & G 개선 · 타 expert도 유효 \\
{}e7 & xss & 0.7522 & 0.7116 & 0.7613 & 0.0148 & 0.0072 & 0.0199 & G 대비 우위 미확인 \\
\bottomrule
\end{longtable}
\endgroup
\end{landscape}
“주로 포함된 클래스”는 context 특화 block의 최다 클래스다. 실제 정답 배정에 사용하지 않는다. 타 expert 최고값은 비교용 요약이며 해당 expert를 test에서 선택한 정책의 점수가 아니다.

\begingroup
\footnotesize
\begin{longtable}{>{\raggedright\arraybackslash}p{\dimexpr 0.12612\linewidth-2\tabcolsep\relax}>{\raggedright\arraybackslash}p{\dimexpr 0.14955\linewidth-2\tabcolsep\relax}>{\raggedright\arraybackslash}p{\dimexpr 0.15737\linewidth-2\tabcolsep\relax}>{\raggedright\arraybackslash}p{\dimexpr 0.14174\linewidth-2\tabcolsep\relax}>{\raggedright\arraybackslash}p{\dimexpr 0.12612\linewidth-2\tabcolsep\relax}>{\raggedright\arraybackslash}p{\dimexpr 0.12612\linewidth-2\tabcolsep\relax}>{\raggedright\arraybackslash}p{\dimexpr 0.17299\linewidth-2\tabcolsep\relax}}
\toprule
Expert & 클래스 & 전면 macro-F1 & Precision & Recall & F1 & 다른 클래스에서 온 FP \\ \midrule
\endfirsthead
Expert & 클래스 & 전면 macro-F1 & Precision & Recall & F1 & 다른 클래스에서 온 FP \\ \midrule
\endhead
{}e1 & injection & 0.6696 & 0.8245 & 0.9659 & 0.8896 & 14,165 \\
{}e2 & scanning & 0.6730 & 0.0910 & 0.9021 & 0.1653 & 54,321 \\
{}e3 & dos & 0.6219 & 0.2984 & 0.9980 & 0.4594 & 52,697 \\
{}e4 & injection & 0.6205 & 0.2411 & 0.9974 & 0.3883 & 216,303 \\
{}e5 & ransomware & 0.6726 & 0.0155 & 0.9661 & 0.0305 & 43,586 \\
{}e6 & benign & 0.6703 & 0.9960 & 0.8096 & 0.8932 & 4,771 \\
{}e7 & xss & 0.6859 & 0.5761 & 0.9996 & 0.7309 & 252,906 \\
\bottomrule
\end{longtable}
\endgroup
Val의 음성 점수에서 FPR\ensuremath{\leq}0.1\%를 만족하는 임계값을 고정한 뒤 test에 그대로 적용했다. 실제 test FPR을 먼저 읽고 그때의 recall을 비교한다. 서로 다른 test FPR의 recall만 비교하면 동일 오탐 조건의 우위라고 말할 수 없다.

\begingroup
\footnotesize
\begin{longtable}{>{\raggedright\arraybackslash}p{\dimexpr 0.13128\linewidth-2\tabcolsep\relax}>{\raggedright\arraybackslash}p{\dimexpr 0.15183\linewidth-2\tabcolsep\relax}>{\raggedright\arraybackslash}p{\dimexpr 0.17237\linewidth-2\tabcolsep\relax}>{\raggedright\arraybackslash}p{\dimexpr 0.20662\linewidth-2\tabcolsep\relax}>{\raggedright\arraybackslash}p{\dimexpr 0.15183\linewidth-2\tabcolsep\relax}>{\raggedright\arraybackslash}p{\dimexpr 0.18607\linewidth-2\tabcolsep\relax}}
\toprule
Expert & 클래스 & G test recall & Expert test recall & G test FPR & Expert test FPR \\ \midrule
\endfirsthead
Expert & 클래스 & G test recall & Expert test recall & G test FPR & Expert test FPR \\ \midrule
\endhead
{}e1 & injection & 0.9515 & 0.9588 & 0.189\% & 0.433\% \\
{}e2 & scanning & 0.2554 & 0.6193 & 0.149\% & 0.004\% \\
{}e3 & dos & 0.9884 & 0.9855 & 0.077\% & 0.018\% \\
{}e4 & injection & 0.9515 & 0.9340 & 0.189\% & 0.364\% \\
{}e5 & ransomware & 0.9140 & 0.9196 & 0.451\% & 0.364\% \\
{}e6 & benign & 0.4114 & 0.7037 & 0.167\% & 0.228\% \\
{}e7 & xss & 0.9377 & 0.8858 & 6.693\% & 5.289\% \\
\bottomrule
\end{longtable}
\endgroup
Test 정답 없이 정한 최근접 centroid 영역에서 비교한다. 상수 대조는 각 영역의 학습용 route 행에서 가장 많은 클래스 하나만 항상 출력한다. 이번 수정에서 추가 계산했다. 영역 내 실제 등장 클래스들의 F1 평균을 표시하며, 같은 행의 세 모델끼리 비교한다.

\begingroup
\footnotesize
\begin{longtable}{>{\raggedright\arraybackslash}p{\dimexpr 0.17793\linewidth-2\tabcolsep\relax}>{\raggedright\arraybackslash}p{\dimexpr 0.14414\linewidth-2\tabcolsep\relax}>{\raggedright\arraybackslash}p{\dimexpr 0.13739\linewidth-2\tabcolsep\relax}>{\raggedright\arraybackslash}p{\dimexpr 0.17117\linewidth-2\tabcolsep\relax}>{\raggedright\arraybackslash}p{\dimexpr 0.20495\linewidth-2\tabcolsep\relax}>{\raggedright\arraybackslash}p{\dimexpr 0.16441\linewidth-2\tabcolsep\relax}}
\toprule
담당 expert /\allowbreak{} 영역 & Test 수 & 등장 클래스 수 & G 영역 macro-F1 & Expert 영역 macro-F1 & 학습 다수 클래스 상수 \\ \midrule
\endfirsthead
담당 expert /\allowbreak{} 영역 & Test 수 & 등장 클래스 수 & G 영역 macro-F1 & Expert 영역 macro-F1 & 학습 다수 클래스 상수 \\ \midrule
\endhead
{}e1 & 76,017 & 6 & 0.3637 & 0.3504 & 0.1547 \\
{}e2 & 10,048 & 9 & 0.2877 & 0.3191 & 0.0612 \\
{}e3 & 47,541 & 9 & 0.1588 & 0.1867 & 0.0707 \\
{}e4 & 40,748 & 10 & 0.1514 & 0.1848 & 0.0079 \\
{}e5 & 8,544 & 3 & 0.0482 & 0.0488 & 0.0470 \\
{}e6 & 1,976,859 & 10 & 0.4793 & 0.4937 & 0.0192 \\
{}e7 & 111,966 & 9 & 0.1095 & 0.1777 & 0.0457 \\
\bottomrule
\end{longtable}
\endgroup
CIC 영역 3·5는 test에 한 클래스만 있어 상수 예측도 1.0000이다. 이런 영역의 높은 점수는 분류 역량의 증거가 되지 않는다. 영역별 등장 클래스가 달라 이 수치를 서로 다른 영역이나 전체 test macro-F1과 직접 비교하지 않는다. Global·모든 타 expert의 영역 accuracy 원본 행렬도 CSV에 보존했다.

\subsection*{2.5. 현재 결과로 내릴 수 있는 판단}
현재 bank 전체가 충분하다는 결론은 내릴 수 없다. 근거가 없는 부분을 다음처럼 구분해야 한다. 이미 식별력 저하가 관측된 대상은 context를 개선하고, 상대 우위가 있는 대상은 고정 운영점과 독립 반복으로 검증한다.

\begingroup
\footnotesize
\begin{longtable}{>{\raggedright\arraybackslash}p{\dimexpr 0.19643\linewidth-2\tabcolsep\relax}>{\raggedright\arraybackslash}p{\dimexpr 0.26786\linewidth-2\tabcolsep\relax}>{\raggedright\arraybackslash}p{\dimexpr 0.26786\linewidth-2\tabcolsep\relax}>{\raggedright\arraybackslash}p{\dimexpr 0.26786\linewidth-2\tabcolsep\relax}}
\toprule
대상 & 현재 근거 & 지금 가능한 판단 & 다음 확인 \\ \midrule
\endfirsthead
대상 & 현재 근거 & 지금 가능한 판단 & 다음 확인 \\ \midrule
\endhead
{}CIC infiltration e6 & AP 0.2178 \ensuremath{\rightarrow} 0.0944; R@0.1\% 0.3007 \ensuremath{\rightarrow} 0.0991 & G 대비 식별력 보완 근거가 없다. 0.9140 정답 배정 F1을 근거로 쓰지 않는다. & K<C로 재생성하고 클래스 비율을 맞춘 무작위 context·어려운 음성 대조 \\
{}ToN scanning e2 & AP 0.3310 \ensuremath{\rightarrow} 0.8312; R@0.1\% 0.1415 \ensuremath{\rightarrow} 0.7462 & 이 holdout에서 G·타 expert 대비 점수 구분 능력의 이득이 있다. & 같은 예산·클래스 비율 대조와 독립 seed/\allowbreak{}시간 분할 \\
{}ToN ransomware e5 & AP 0.0477 \ensuremath{\rightarrow} 0.2267; R@0.1\% 0.0310 \ensuremath{\rightarrow} 0.6996 & 순위 이득은 있지만 val 0.1\% 임계값의 test FPR은 0.364\%로 목표를 넘는다. & 임계값 전이 검증·시간 변화 보정 후 새 test 재평가 \\
{}CIC bank & 고정 영역 0.7477; G/\allowbreak{}S/\allowbreak{}S+\allowbreak{}V 0.8107 & 기존 bank·고정 영역의 결합은 G보다 낮고 실제 정책은 개입하지 않는다. & 새 K 조건의 bank 재학습 및 정답 없는 validation 비교 \\
\bottomrule
\end{longtable}
\endgroup
“충분”을 최종 판정하려면 목표 클래스·허용 FPR·최소 recall·허용 global 성능 저하를 평가 전에 정해야 한다. 지금의 AP 또는 R@0.1\%에 임의의 합격선을 사후 부여하지 않는다. 다음 실험은 그 판정에 필요한 원인·전이·재현성 근거를 채우기 위한 것이다.

\subsection*{2.6. 추가로 필요한 실험과 각각의 판정 기준}
\begingroup
\footnotesize
\begin{longtable}{>{\raggedright\arraybackslash}p{\dimexpr 0.23694\linewidth-2\tabcolsep\relax}>{\raggedright\arraybackslash}p{\dimexpr 0.26679\linewidth-2\tabcolsep\relax}>{\raggedright\arraybackslash}p{\dimexpr 0.26679\linewidth-2\tabcolsep\relax}>{\raggedright\arraybackslash}p{\dimexpr 0.22948\linewidth-2\tabcolsep\relax}}
\toprule
우선순위 /\allowbreak{} 실험 & 무엇을 고정하고 비교하는가 & 무엇이 확인돼야 하는가 & 상태 \\ \midrule
\endfirsthead
우선순위 /\allowbreak{} 실험 & 무엇을 고정하고 비교하는가 & 무엇이 확인돼야 하는가 & 상태 \\ \midrule
\endhead
{}0. 정답 없는 배정 +\allowbreak{} 상수 expert 대조 & 동일 cache·고정 영역에서 원래 expert와 학습 다수 클래스 상수 예측 비교 & 높은 점수가 영역의 단일 클래스 구성만으로 설명되는지 확인 & 이번 수정에서 계산 완료 \\
{}1. K<C bank 재생성과 context 선택 대조 & 같은 K·anchor·총 context 행 수·클래스별 행 수에서 현재 실패 패턴 선택 vs 무작위 선택; 별도 동일 크기 균형 context 대조 & 현재 방식이 클래스 비율을 맞춘 무작위보다 AP·낮은 FPR recall에서 반복적으로 높아야 생성 방식의 효과를 지지 & 추가 학습 필요 \\
{}2. Global 보정만 한 대조 & 새 expert 없이 global의 prior/\allowbreak{}β·온도·클래스 임계값만 validation에서 보정; expert도 같은 선택 절차 사용 & expert 이득이 global의 운영점 변경만으로 설명되는지 구별 & 추가 대조 필요 \\
{}3. 고정 운영점의 시간 전이 & 목표 FPR 0.1\%·1\%를 먼저 정하고 val에서 임계값 고정  \ensuremath{\rightarrow}  이후 test에서 recall·실제 FPR·FP 수 평가 & 허용 FPR을 지키면서 G보다 높은 recall이 유지돼야 그 운영점의 이득을 지지 & 현재 holdout 계산 완료; 새 시간 분할 필요 \\
{}4. 독립 학습 반복·그룹 불확실성 & 최소 3개의 실제 context 추첨·학습 seed; 새 시간 holdout; 같은 벡터·시간 블록 단위 paired bootstrap & 평균 이득과 불확실성 구간을 함께 제시하고 특정 seed·중복 행에만 의존하는지 확인 & 추가 학습·새 평가 필요 \\
{}선택. Bank 내 개별 expert 제거 & 해당 expert를 제거하고 같은 validation 절차로 비교 정책 재선택 & bank에 꼭 필요한 expert라는 주장에만 필요; 단독 식별력 검증을 대신하지 않는다 & 고유 시스템 기여를 주장할 때 실행 \\
\bottomrule
\end{longtable}
\endgroup
생성 품질을 주장하기 위한 우선 조합은 1·2·3·4번이다. 새 K 조건을 만족하는 bank를 만들고, 동일 예산 대조에서 식별력 이득이 남는지와 새 데이터에서도 정한 오탐 한도를 지키는지를 확인한다. 현재 이미 계산된 0번과 기존 확률 진단을 새 학습 결과로 세지 않는다.

기존 EXP52의 행 선택·anchor 대조와 EXP53–54의 context 구성·prior 대조는 이미 수행된 부분 근거다. 이를 폐기하거나 미실행으로 세지 않는다. 여기서 추가로 요구하는 것은 K<C의 동일 bank 예산·클래스 구성 통제, 같은 저오탐 지표, 독립 반복·새 시간 평가를 갖춘 비교다.

기존 correctness oracle과 정답 클래스 배정은 주 품질 비교에서 모두 제외했다. 전자는 정오를 보고 유리한 답을 고르고, 후자는 정답 클래스로 양성과 음성을 분리한다. 과거 수치는 재현·차이 설명용으로만 보존한다. 0911 원본 데이터 감사의 과거 참조값도 별도로 구분한다.

이번 수정의 2번·상수 대조·예측 출처 추적: tabpfn/\allowbreak{}scripts/\allowbreak{}exp56\_\allowbreak{}label\_\allowbreak{}free\_\allowbreak{}revision.py . 새 평가 기본 규칙: tabpfn/\allowbreak{}scripts/\allowbreak{}exp56\_\allowbreak{}five\_\allowbreak{}condition\_\allowbreak{}evaluation.py . 결과: tabpfn/\allowbreak{}results/\allowbreak{}exp56\_\allowbreak{}label\_\allowbreak{}free\_\allowbreak{}20260922 . AP·FPR 진단은 동일 cache의 기존 CSV를 hash와 함께 재사용했다. 기존 조건 1·3·4·5의 점수와 S/\allowbreak{}S+\allowbreak{}V 예측 배열은 원본과 일치한다. K 제약: tabpfn/\allowbreak{}scripts/\allowbreak{}nfv3\_\allowbreak{}v3\_\allowbreak{}exp31\_\allowbreak{}c0alloc.py . 개수 제약만 구현했으며 K 선택 목적함수를 새로 최적화하거나 기존 bank를 재학습하지 않았다.

\subsection*{0911 원본 데이터의 과거 정의 — 위 모순 제거 평가와 구분}
09-22 정정. 기존 oracle은 정답 여부를 보고 유리한 예측만 채택하므로 expert가 담당 영역에서 잘 생성됐는지 평가하는 기준으로 쓰지 않는다. 정답 클래스를 배정에 제공하는 oracle도 양성·음성을 미리 분리하므로 주 평가에서 제외했다. 현재 2번은 정답 없는 고정 영역 배정(S/\allowbreak{}V 없음)이며 oracle로 부르지 않는다. 구체적 대조와 재해석은 통합본 목차의 0918 수정 보고서 에 제시한다.

\begingroup
\footnotesize
\begin{longtable}{>{\raggedright\arraybackslash}p{\dimexpr 0.31312\linewidth-2\tabcolsep\relax}>{\raggedright\arraybackslash}p{\dimexpr 0.20421\linewidth-2\tabcolsep\relax}>{\raggedright\arraybackslash}p{\dimexpr 0.25866\linewidth-2\tabcolsep\relax}>{\raggedright\arraybackslash}p{\dimexpr 0.22401\linewidth-2\tabcolsep\relax}}
\toprule
기존 항목 & Bank & Test 라벨 개입 & 새 명칭·용도 \\ \midrule
\endfirsthead
기존 항목 & Bank & Test 라벨 개입 & 새 명칭·용도 \\ \midrule
\endhead
{}CIC row\_\allowbreak{}level\_\allowbreak{}oracle\_\allowbreak{}bank · 0.849159 & Global+\allowbreak{}expert 3개 & 샘플별 정답을 낸 후보 선택 & Bank 정답 포함률의 합성 F1 \\
{}CIC realistic\_\allowbreak{}oracle\_\allowbreak{}bank · 0.803623 & 같은 bank & 충돌은 test 다수 라벨을 목표로 선택 & Bank+\allowbreak{}다수 라벨 정책 참조값 \\
{}CIC realistic\_\allowbreak{}oracle\_\allowbreak{}perfect · 0.855907 & 없음 & 비충돌에는 y, 충돌에는 test 다수 라벨 대입 & 데이터 다수 라벨 참조값 \\
{}ToN·BoT·UNSW /\allowbreak{} EXP41·42·44 참조 & 없음 & Test 벡터별 다수 라벨 대입 & 데이터 진단; expert 성능 아님 \\
\bottomrule
\end{longtable}
\endgroup
다수 라벨은 accuracy 기준의 배정이며 macro-F1 상한과 같지 않다 . 클래스별로 별도 F1 최대화를 한 평균은 느슨한 macro 상계다(CIC bank 제약에서는 0.804844). 같은 벡터에는 같은 답을 내는 제약과 “train의 다수 라벨을 반드시 따라야 한다”는 가정도 다르다. 후자는 특정 조회 정책의 가정이며 모든 학습기의 제약이 아니다. Bank 없는 참조값이 모순 제거 후 1.0이 되어도 학습·일반화 가능성을 보장하지 않는다.

0911 원본 CIC EXP39와 0918 모순 제거 CIC EXP47·ToN EXP48은 test·bank가 다르다. 날짜 간 점수 차이를 같은 oracle의 개선량으로 비교하지 않는다. 과거 0821·0902 수치도 bank와 선택 규칙을 각각 명시해야 한다.

\begingroup
\footnotesize
\begin{longtable}{>{\raggedright\arraybackslash}p{\dimexpr 0.16441\linewidth-2\tabcolsep\relax}>{\raggedright\arraybackslash}p{\dimexpr 0.13063\linewidth-2\tabcolsep\relax}>{\raggedright\arraybackslash}p{\dimexpr 0.17117\linewidth-2\tabcolsep\relax}>{\raggedright\arraybackslash}p{\dimexpr 0.21847\linewidth-2\tabcolsep\relax}>{\raggedright\arraybackslash}p{\dimexpr 0.15090\linewidth-2\tabcolsep\relax}>{\raggedright\arraybackslash}p{\dimexpr 0.16441\linewidth-2\tabcolsep\relax}}
\toprule
클래스 & 실제 F1 & Bank 정답 선택 F1 & 같은 bank +\allowbreak{} test 다수 라벨 & 클래스별 F1 상한 & Train 라벨 충돌률 \\ \midrule
\endfirsthead
클래스 & 실제 F1 & Bank 정답 선택 F1 & 같은 bank +\allowbreak{} test 다수 라벨 & 클래스별 F1 상한 & Train 라벨 충돌률 \\ \midrule
\endhead
{}benign & 0.995 & 0.999 & 0.997 & 0.997 & 0.2\% \\
{}bot & 0.998 & 0.999 & 0.999 & 0.999 & 0.0\% \\
{}brute\_\allowbreak{}force & 0.916 & 0.916 & 0.916 & 0.916 & 0.0\% \\
{}ddos & 0.996 & 1.000 & 1.000 & 1.000 & 0.0\% \\
{}dos & 0.781 & 0.782 & 0.782 & 0.782 & 34.8\% \\
{}infiltration & 0.400 & 0.970 & 0.666 & 0.672 & 49.5\% \\
{}web\_\allowbreak{}attacks & 0.248 & 0.278 & 0.266 & 0.269 & 30.3\% \\
\bottomrule
\end{longtable}
\endgroup
CIC EXP39 원본 test의 global F1은 0.762086, bank의 행별 정답 선택 F1은 0.849159, bank에 test 다수 라벨 선택 규칙을 적용한 참조값은 0.803623다. Bank 없이 test 라벨을 직접 대입하되 모순 벡터에만 다수 라벨을 쓰면 0.855907로, 별개 진단이다. Bank·동일 벡터 제약의 클래스별 max-F1 평균 0.804844도 다수 라벨 참조값과 구분한다. EXP39의 교정 가능한 infiltration 25,423행 중 17,564행이 benign 657,486행과 벡터를 공유한다. 이 수치들은 expert의 담당 영역 성능이나 실현 가능한 라우팅 이득을 직접 측정하지 않는다.

\begingroup
\footnotesize
\begin{longtable}{>{\raggedright\arraybackslash}p{\dimexpr 0.15476\linewidth-2\tabcolsep\relax}>{\raggedright\arraybackslash}p{\dimexpr 0.13333\linewidth-2\tabcolsep\relax}>{\raggedright\arraybackslash}p{\dimexpr 0.17619\linewidth-2\tabcolsep\relax}>{\raggedright\arraybackslash}p{\dimexpr 0.21190\linewidth-2\tabcolsep\relax}>{\raggedright\arraybackslash}p{\dimexpr 0.15476\linewidth-2\tabcolsep\relax}>{\raggedright\arraybackslash}p{\dimexpr 0.16905\linewidth-2\tabcolsep\relax}}
\toprule
클래스 & 실제 F1 & Bank 정답 선택 F1 & bank 없이 test 다수 라벨 & 클래스별 F1 상한 & Train 라벨 충돌률 \\ \midrule
\endfirsthead
클래스 & 실제 F1 & Bank 정답 선택 F1 & bank 없이 test 다수 라벨 & 클래스별 F1 상한 & Train 라벨 충돌률 \\ \midrule
\endhead
{}backdoor & 0.998 & — & 0.999 & 0.999 & 0.1\% \\
{}benign & 0.914 & — & 0.986 & 0.986 & 3.3\% \\
{}ddos & 0.987 & — & 0.994 & 0.994 & 1.6\% \\
{}dos & 0.817 & — & 0.902 & 0.903 & 28.1\% \\
{}injection & 0.941 & — & 0.973 & 0.973 & 5.5\% \\
{}mitm & 0.162 & — & 0.987 & 0.987 & 3.2\% \\
{}password & 0.842 & — & 0.920 & 0.922 & 0.1\% \\
{}ransomware & 0.150 & — & 0.948 & 0.949 & 1.4\% \\
{}scanning & 0.001 & — & 1.000 & 1.000 & 97.8\% \\
{}xss & 0.876 & — & 0.964 & 0.964 & 0.6\% \\
\bottomrule
\end{longtable}
\endgroup
이 절의 ToN 다수 라벨 참조값에는 expert bank가 없다. Scanning은 test 내부의 같은 벡터 라벨 제약만 보면 참조 F1이 1.00지만, test 행 중 98\%는 같은 벡터의 train 다수 라벨과 충돌한다. 실제 full-XGB scanning F1은 0.001이다. Mitm·ransomware의 참조 F1 0.99·0.95와 실제 F1 0.16·0.15의 격차는 데이터 진단이며, 이만큼을 현재 expert나 라우터가 회수할 수 있다는 의미가 아니다. 오탐 원인과 시간 이동을 분리해 평가한다.

\begingroup
\footnotesize
\begin{longtable}{>{\raggedright\arraybackslash}p{\dimexpr 0.17793\linewidth-2\tabcolsep\relax}>{\raggedright\arraybackslash}p{\dimexpr 0.13063\linewidth-2\tabcolsep\relax}>{\raggedright\arraybackslash}p{\dimexpr 0.17117\linewidth-2\tabcolsep\relax}>{\raggedright\arraybackslash}p{\dimexpr 0.20495\linewidth-2\tabcolsep\relax}>{\raggedright\arraybackslash}p{\dimexpr 0.15090\linewidth-2\tabcolsep\relax}>{\raggedright\arraybackslash}p{\dimexpr 0.16441\linewidth-2\tabcolsep\relax}}
\toprule
클래스 & 실제 F1 & Bank 정답 선택 F1 & bank 없이 test 다수 라벨 & 클래스별 F1 상한 & Train 라벨 충돌률 \\ \midrule
\endfirsthead
클래스 & 실제 F1 & Bank 정답 선택 F1 & bank 없이 test 다수 라벨 & 클래스별 F1 상한 & Train 라벨 충돌률 \\ \midrule
\endhead
{}benign & 0.826 & — & 0.854 & 0.854 & 25.1\% \\
{}ddos & 0.031 & — & 0.999 & 0.999 & 0.0\% \\
{}dos & 0.677 & — & 0.998 & 0.998 & 1.6\% \\
{}reconnaissance & 0.884 & — & 0.998 & 0.998 & 0.0\% \\
{}theft & 0.387 & — & 0.982 & 0.982 & 74.9\% \\
\bottomrule
\end{longtable}
\endgroup
Theft의 test 다수 라벨 참조 F1은 0.98지만, 같은 벡터가 train에서는 recon으로 라벨된 행의 비율이 75\%다. Full-XGB theft F1은 0.387이다. 높은 정답 대입 참조값을 근거로 expert 품질이나 학습 가능성을 판단하지 않는다. 시간 구간별 라벨·특징 분포를 함께 확인한다.

\begingroup
\footnotesize
\begin{longtable}{>{\raggedright\arraybackslash}p{\dimexpr 0.17793\linewidth-2\tabcolsep\relax}>{\raggedright\arraybackslash}p{\dimexpr 0.13063\linewidth-2\tabcolsep\relax}>{\raggedright\arraybackslash}p{\dimexpr 0.17117\linewidth-2\tabcolsep\relax}>{\raggedright\arraybackslash}p{\dimexpr 0.20495\linewidth-2\tabcolsep\relax}>{\raggedright\arraybackslash}p{\dimexpr 0.15090\linewidth-2\tabcolsep\relax}>{\raggedright\arraybackslash}p{\dimexpr 0.16441\linewidth-2\tabcolsep\relax}}
\toprule
클래스 & 실제 F1 & Bank 정답 선택 F1 & bank 없이 test 다수 라벨 & 클래스별 F1 상한 & Train 라벨 충돌률 \\ \midrule
\endfirsthead
클래스 & 실제 F1 & Bank 정답 선택 F1 & bank 없이 test 다수 라벨 & 클래스별 F1 상한 & Train 라벨 충돌률 \\ \midrule
\endhead
{}analysis & 0.282 & — & 1.000 & 1.000 & 0.0\% \\
{}backdoor & 0.008 & — & 0.797 & 0.798 & 39.5\% \\
{}benign & 1.000 & — & 1.000 & 1.000 & 0.0\% \\
{}dos & 0.362 & — & 0.981 & 0.982 & 4.1\% \\
{}exploits & 0.726 & — & 0.980 & 0.980 & 3.8\% \\
{}fuzzers & 0.752 & — & 0.964 & 0.965 & 8.1\% \\
{}generic & 0.860 & — & 0.966 & 0.966 & 1.1\% \\
{}reconnaissance & 0.677 & — & 0.977 & 0.977 & 3.4\% \\
{}shellcode & 0.736 & — & 0.952 & 0.953 & 6.3\% \\
{}worms & 0.708 & — & 0.968 & 0.968 & 6.2\% \\
\bottomrule
\end{longtable}
\endgroup
UNSW의 참조값도 expert bank 없이 test 벡터별 다수 라벨을 대입한 값이다(macro 0.959). 실제 XGB macro는 0.611이다. 정확일치 라벨 모순이 적더라도 근사 중복과 클래스 간 혼동은 남을 수 있으므로 참조값과의 차이를 학습 가능한 여지로 단정하지 않는다.

모순·중복·누수를 다른 연구는 어떻게 처리했는가. 아래는 링크 정리이며 보충 예정 이다.

전체 목록과 확인 수준(본문 확인 /\allowbreak{} 초록만 /\allowbreak{} 미확인)은 docs/\allowbreak{}lablog/\allowbreak{}reports/\allowbreak{}20260911.tex 에 있다. NetFlow v2에 대한 중복·모순 감사는 Raskovalov 외의 한 문장이 전부이고, v3에 대해서는 전무하다 — Sarhan 외의 표준 feature set 논문과 Luay 외의 temporal 논문 어디에도 duplicate·redundant·identical·conflict가 나오지 않는다.

\subsection*{4. 중복과 모순을 분리해서 처리하면 — EXP41·EXP42}
§1의 직교성 위에 처리를 놓으면 네 칸이 나온다. 중복을 지우는 처리는 "중복을 지운 효과"와 "모순을 없앤 효과"를 한꺼번에 섞으므로, 둘을 떼어 보려면 중복을 남긴 채 모순만 없애는 칸이 필요하다.

\begingroup
\footnotesize
\begin{longtable}{>{\raggedright\arraybackslash}p{\dimexpr 0.20884\linewidth-2\tabcolsep\relax}>{\raggedright\arraybackslash}p{\dimexpr 0.39558\linewidth-2\tabcolsep\relax}>{\raggedright\arraybackslash}p{\dimexpr 0.39558\linewidth-2\tabcolsep\relax}}
\toprule
 & 모순 유지 & 모순 제거 \\ \midrule
\endfirsthead
 & 모순 유지 & 모순 제거 \\ \midrule
\endhead
{}중복 유지 & 원본 train 12,069,313행 & 중복 유지 +\allowbreak{} 모순은 다수 라벨 모순 벡터의 행을 그 벡터의 다수 라벨로 바꾼다. 행을 하나도 지우지 않으므로 train·test 행이 원본과 동일하다. 재라벨 210,728행. train 12,069,313행 (EXP44) 중복 유지 +\allowbreak{} 모순 벡터 삭제 모순 벡터의 행만 제거하고 나머지 중복은 남긴다 (EXP44) \\
{}중복 제거 & 중복 제거 벡터·라벨 쌍당 1행. 모순 벡터는 라벨별로 1행씩 남으므로 모순은 그대로다 . train 5,295,822행 & 중복 제거 +\allowbreak{} 모순은 다수 라벨 벡터당 1행, 라벨은 그 벡터의 다수 라벨. train 5,283,411행 중복 제거 +\allowbreak{} 모순 벡터 삭제 모순 벡터를 통째로 제거. train 5,271,045행 \\
\bottomrule
\end{longtable}
\endgroup
각 처리를 전 행에 적용한 뒤 같은 scenario별 시간순 60/\allowbreak{}20/\allowbreak{}20 규칙으로 다시 나눈다. 여섯 처리는 서로 다른 벤치마크가 되므로 먼저 클래스별로 train·test가 몇 행이 됐는지 를 보이고, 그 위에서 XGBoost(train 전체 /\allowbreak{} 같은 C0) · TabPFN-v3(무작위 컨텍스트) · ours(구성된 C0) · SOTA 3종을 클래스별로 비교한다(EXP42; 중복 유지 +\allowbreak{} 모순 제거 두 arm은 EXP44이며 LoCalPFN은 arm당 \textasciitilde{}3.4 h라 제외). 각 arm의 test로 평가하고, bank 없는 test 다수 라벨 참조값을 별도 데이터 진단으로 둔다. 여기서 ours =\allowbreak{} 구성된 C0 위의 frozen TabPFN-v3 다 — 조건부 라우팅 층은 09-09의 EXP38·39에서 6조건 모두 global을 유지해 기여가 0이었으므로(Δglobal =\allowbreak{} 0), 남은 기여는 컨텍스트 구성뿐이고 그것을 그대로 비교 대상으로 둔다. 이전 기록: EXP16(0821, train만 중복 제거)은 global까지 붕괴(0.701 \ensuremath{\rightarrow} 0.599), EXP18(0821, 분할 전 클래스 내 dedup)은 macro 0.7547, exp32/\allowbreak{}33(0902–03)은 음성.

\begin{landscape}
\begingroup
\scriptsize
\begin{longtable}{>{\raggedright\arraybackslash}p{\dimexpr 0.06802\linewidth-2\tabcolsep\relax}>{\raggedright\arraybackslash}p{\dimexpr 0.06309\linewidth-2\tabcolsep\relax}>{\raggedright\arraybackslash}p{\dimexpr 0.08033\linewidth-2\tabcolsep\relax}>{\raggedright\arraybackslash}p{\dimexpr 0.08280\linewidth-2\tabcolsep\relax}>{\raggedright\arraybackslash}p{\dimexpr 0.08033\linewidth-2\tabcolsep\relax}>{\raggedright\arraybackslash}p{\dimexpr 0.08033\linewidth-2\tabcolsep\relax}>{\raggedright\arraybackslash}p{\dimexpr 0.07787\linewidth-2\tabcolsep\relax}>{\raggedright\arraybackslash}p{\dimexpr 0.07787\linewidth-2\tabcolsep\relax}>{\raggedright\arraybackslash}p{\dimexpr 0.07787\linewidth-2\tabcolsep\relax}>{\raggedright\arraybackslash}p{\dimexpr 0.07787\linewidth-2\tabcolsep\relax}>{\raggedright\arraybackslash}p{\dimexpr 0.07787\linewidth-2\tabcolsep\relax}>{\raggedright\arraybackslash}p{\dimexpr 0.07787\linewidth-2\tabcolsep\relax}>{\raggedright\arraybackslash}p{\dimexpr 0.07787\linewidth-2\tabcolsep\relax}}
\toprule
클래스 & 원본 & 중복 유지 +\allowbreak{} 모순은 다수 라벨 & 중복 유지 +\allowbreak{} 모순 벡터 삭제 & 중복 제거 & 중복 제거 +\allowbreak{} 모순은 다수 라벨 & 중복 제거 +\allowbreak{} 모순 벡터 삭제 &  &  &  &  &  &  \\ \midrule
\endfirsthead
클래스 & 원본 & 중복 유지 +\allowbreak{} 모순은 다수 라벨 & 중복 유지 +\allowbreak{} 모순 벡터 삭제 & 중복 제거 & 중복 제거 +\allowbreak{} 모순은 다수 라벨 & 중복 제거 +\allowbreak{} 모순 벡터 삭제 &  &  &  &  &  &  \\ \midrule
\endhead
{}train & test & train & test & train & test & train & test & train & test & train & test &  \\
{}benign & 10,508,775 & 3,502,926 & 10,564,482 100.5\% & 3,520,680 100.5\% & 7,678,996 73.1\% & 2,559,666 73.1\% & 4,823,483 45.9\% & 1,607,828 45.9\% & 4,821,345 45.9\% & 1,607,116 45.9\% & 4,811,120 45.8\% & 1,603,707 45.8\% \\
{}bot & 124,621 & 41,541 & 124,621 100.0\% & 41,541 100.0\% & 124,621 100.0\% & 41,541 100.0\% & 1,682 1.3\% & 561 1.4\% & 1,682 1.3\% & 561 1.4\% & 1,682 1.3\% & 561 1.4\% \\
{}brute\_\allowbreak{}force & 345,116 & 115,039 & 408,446 118.4\% & 136,149 118.3\% & 113,048 32.8\% & 37,683 32.8\% & 56,388 16.3\% & 18,796 16.3\% & 56,388 16.3\% & 18,796 16.3\% & 56,387 16.3\% & 18,796 16.3\% \\
{}ddos & 794,609 & 264,871 & 795,216 100.1\% & 265,008 100.1\% & 789,614 99.4\% & 263,206 99.4\% & 306,965 38.6\% & 102,323 38.6\% & 306,965 38.6\% & 102,323 38.6\% & 306,964 38.6\% & 102,323 38.6\% \\
{}dos & 181,779 & 60,594 & 118,449 65.2\% & 39,483 65.2\% & 118,449 65.2\% & 39,483 65.2\% & 78,596 43.2\% & 26,199 43.2\% & 78,594 43.2\% & 26,199 43.2\% & 78,594 43.2\% & 26,199 43.2\% \\
{}infiltration & 112,891 & 37,631 & 57,564 51.0\% & 20,086 53.4\% & 28,056 24.9\% & 9,353 24.9\% & 28,118 24.9\% & 9,373 24.9\% & 17,995 15.9\% & 5,999 15.9\% & 15,889 14.1\% & 5,297 14.1\% \\
{}web\_\allowbreak{}attacks & 1,522 & 508 & 535 35.2\% & 163 32.0\% & 461 30.3\% & 156 30.6\% & 590 38.8\% & 197 38.7\% & 442 29.0\% & 149 29.3\% & 409 26.9\% & 138 27.1\% \\
{}합계 & 12,069,313 & 4,023,110 & 12,069,313 100.0\% & 4,023,110 100.0\% & 8,853,245 73.4\% & 2,951,088 73.4\% & 5,295,822 43.9\% & 1,765,277 43.9\% & 5,283,411 43.8\% & 1,761,143 43.8\% & 5,271,045 43.7\% & 1,757,021 43.7\% \\
\bottomrule
\end{longtable}
\endgroup
\end{landscape}
각 처리를 전 행에 적용한 뒤 같은 scenario별 시간순 60/\allowbreak{}20/\allowbreak{}20 규칙으로 다시 나눈 결과다(행이 5개 미만인 시나리오는 train 전용). 작은 숫자는 파이프라인 원본 split의 같은 칸 대비 남은 비율이다. val(20\%)은 표에서 생략했다. 원본 arm의 test는 재구현 split이라 파이프라인과 4행 차이가 난다.

\begingroup
\footnotesize
\begin{longtable}{>{\raggedright\arraybackslash}p{\dimexpr 0.23245\linewidth-2\tabcolsep\relax}>{\raggedright\arraybackslash}p{\dimexpr 0.10109\linewidth-2\tabcolsep\relax}>{\raggedright\arraybackslash}p{\dimexpr 0.14346\linewidth-2\tabcolsep\relax}>{\raggedright\arraybackslash}p{\dimexpr 0.13923\linewidth-2\tabcolsep\relax}>{\raggedright\arraybackslash}p{\dimexpr 0.10109\linewidth-2\tabcolsep\relax}>{\raggedright\arraybackslash}p{\dimexpr 0.14346\linewidth-2\tabcolsep\relax}>{\raggedright\arraybackslash}p{\dimexpr 0.13923\linewidth-2\tabcolsep\relax}}
\toprule
방법 (macro-F1) & 원본 & 중복 유지 +\allowbreak{} 모순은 다수 라벨 & 중복 유지 +\allowbreak{} 모순 벡터 삭제 & 중복 제거 & 중복 제거 +\allowbreak{} 모순은 다수 라벨 & 중복 제거 +\allowbreak{} 모순 벡터 삭제 \\ \midrule
\endfirsthead
방법 (macro-F1) & 원본 & 중복 유지 +\allowbreak{} 모순은 다수 라벨 & 중복 유지 +\allowbreak{} 모순 벡터 삭제 & 중복 제거 & 중복 제거 +\allowbreak{} 모순은 다수 라벨 & 중복 제거 +\allowbreak{} 모순 벡터 삭제 \\ \midrule
\endhead
{}다수 라벨 참조값 bank 없음 · test 벡터별 다수 라벨 직접 대입 & 0.8564 & 1.0000 & 1.0000 & 0.9974 & 1.0000 & 1.0000 \\
{}XGBoost train 풀 전체 (300 trees · depth 8) & 0.7548 & 0.8467 & 0.7968 & 0.7613 & 0.7759 & 0.7763 \\
{}XGBoost 같은 C0 100k행 (정보 예산 짝지음) & 0.7525 & 0.7774 & 0.7921 & 0.8043 & 0.7956 & 0.7919 \\
{}TabPFN-v3 무작위 100k 컨텍스트 (비율 보존) & 0.7282 & 0.8055 & 0.7590 & 0.7411 & 0.7321 & 0.7541 \\
{}ours 구성된 C0 100k (benign 0.75 · 공격 균등 · n\_\allowbreak{}est 4) & 0.7616 & 0.8140 & 0.7923 & 0.8055 & 0.7977 & 0.8039 \\
{}DistPFN ours(n\_\allowbreak{}est 4) 확률에 prior 보정 & 0.6681 & 0.8106 & 0.7931 & 0.7590 & 0.6150 & 0.6041 \\
{}BoostPFN TabPFN-v1 약학습기 T=\allowbreak{}50 × 500행 & 0.7368 & 0.7708 & 0.7194 & 0.7404 & 0.6884 & 0.6636 \\
{}LoCalPFN TabPFN-kNN: 행마다 국소 컨텍스트 1,000행 & 0.7121 & — & — & 0.7690 & 0.7588 & 0.7402 \\
\bottomrule
\end{longtable}
\endgroup
열마다 test가 다르므로 가로 비교는 "그 벤치마크에서의 점수"다. 같은 열 안의 세로 비교만 방법 간 비교다. 각 열에서 oracle을 제외한 최고값을 굵게 표시했다. seed 42, XGBoost 300 trees · depth 8. 컨텍스트를 쓰는 다섯 방법은 모두 같은 크기(100k행)를 받는다 — XGBoost C0 · ours · BoostPFN · LoCalPFN은 구성된 C0를, TabPFN-v3는 같은 풀에서 비율을 보존해 무작위로 뽑은 100k를 받는다. 그 차이가 컨텍스트 구성의 몫이다. DistPFN은 ours의 확률에 사후 prior 보정을 건 것인데, exp37 기록(원본 벤치마크)에서는 n\_\allowbreak{}estimators 4·8·16·32 모두 중립\textasciitilde{}상승(0.7617 \ensuremath{\rightarrow} 0.7609, 0.7635 \ensuremath{\rightarrow} 0.7633, 0.7631 \ensuremath{\rightarrow} 0.7691, 0.7693 \ensuremath{\rightarrow} 0.7692)이었던 것과 달리 EXP42 재구현에서는 원본 arm 0.6681로 내려간다 — n\_\allowbreak{}estimators로는 설명되지 않으며 exp37 플러그인과 EXP42 재구현의 차이는 미해결이다(0915 정정). C0가 꼬리를 일부러 과대표집한 컨텍스트라 "prior shift 보정"이 그 구성을 되돌리는 방향으로 작동하는 것은 두 구현에 공통이다.

재라벨 arm에서는 train 풀 전체를 쓴 XGBoost(0.8467)가 ours(0.8140)와 무작위 컨텍스트 TabPFN-v3(0.8055)를 앞선다. 원본에서는 ours가 앞섰으므로 모순을 없애면 순위가 뒤집힌다 — 12M행을 통째로 외우는 쪽이 유리해지고 컨텍스트 구성의 이점이 줄어든다. 벡터 삭제 arm에서는 XGBoost 0.7968 · ours 0.7923 · XGB-c0 0.7921가 사실상 동률이다. 즉 ours의 우위는 모순이 남아 있을 때만 나타난다 — 우리 주장에 불리하지만 이것이 데이터가 말하는 바다. 그리고 재라벨은 행을 하나도 지우지 않는데도 test web 509 \ensuremath{\rightarrow} 163행(32.0\%), infiltration 37,631 \ensuremath{\rightarrow} 20,086행(53.4\%)이 다수 라벨(benign)로 흡수되고 dos의 slowhttptest는 brute\_\allowbreak{}force로 넘어간다. 두 arm 모두 macro 상승분의 상당 부분은 모델이 좋아져서가 아니라 어려운 행이 사라지거나 쉬운 클래스로 옮겨간 몫이다. 다수 라벨은 train+\allowbreak{}test 전체에서 계산하므로 재라벨 arm과 중복 제거+\allowbreak{}다수 라벨 arm은 모두 Arp et al.(USENIX Sec'22)의 selective snooping에 해당하는 진단용 벤치마크이지 배포 가능한 전처리가 아니다. 결론: 모순은 지워도 고쳐도 꼬리를 잃는다. headline은 원본 벤치마크로 두고, 모순은 §1의 3분해와 클래스별 다수 라벨 참조값을 옆에 두는 방식으로 보고한다(§7). 데이터 감사 실험은 여기서 마친다(09-15).

\begin{landscape}
\begingroup
\scriptsize
\begin{longtable}{>{\raggedright\arraybackslash}p{\dimexpr 0.23627\linewidth-2\tabcolsep\relax}>{\raggedright\arraybackslash}p{\dimexpr 0.08922\linewidth-2\tabcolsep\relax}>{\raggedright\arraybackslash}p{\dimexpr 0.08431\linewidth-2\tabcolsep\relax}>{\raggedright\arraybackslash}p{\dimexpr 0.08431\linewidth-2\tabcolsep\relax}>{\raggedright\arraybackslash}p{\dimexpr 0.08431\linewidth-2\tabcolsep\relax}>{\raggedright\arraybackslash}p{\dimexpr 0.08431\linewidth-2\tabcolsep\relax}>{\raggedright\arraybackslash}p{\dimexpr 0.08431\linewidth-2\tabcolsep\relax}>{\raggedright\arraybackslash}p{\dimexpr 0.08431\linewidth-2\tabcolsep\relax}>{\raggedright\arraybackslash}p{\dimexpr 0.08431\linewidth-2\tabcolsep\relax}>{\raggedright\arraybackslash}p{\dimexpr 0.08431\linewidth-2\tabcolsep\relax}}
\toprule
방법 & macro-F1 & tail F1 & benign & bot & brute & ddos & dos & inf & web \\ \midrule
\endfirsthead
방법 & macro-F1 & tail F1 & benign & bot & brute & ddos & dos & inf & web \\ \midrule
\endhead
{}다수 라벨 참조값 bank 없음 · test 벡터별 다수 라벨 직접 대입 & 0.8564 & 0.7641 & 0.997 & 1.000 & 0.916 & 1.000 & 0.789 & 0.706 & 0.586 \\
{}XGBoost train 풀 전체 (300 trees · depth 8) & 0.7548 & 0.5331 & 0.994 & 0.998 & 0.916 & 0.992 & 0.782 & 0.429 & 0.172 \\
{}XGBoost 같은 C0 100k행 (정보 예산 짝지음) & 0.7525 & 0.5259 & 0.992 & 0.997 & 0.916 & 0.999 & 0.782 & 0.337 & 0.244 \\
{}TabPFN-v3 무작위 100k 컨텍스트 (비율 보존) & 0.7282 & 0.4702 & 0.995 & 0.998 & 0.916 & 0.994 & 0.782 & 0.413 & 0.000 \\
{}ours 구성된 C0 100k (benign 0.75 · 공격 균등 · n\_\allowbreak{}est 4) & 0.7616 & 0.5478 & 0.995 & 0.998 & 0.916 & 0.996 & 0.781 & 0.400 & 0.245 \\
{}DistPFN ours(n\_\allowbreak{}est 4) 확률에 prior 보정 & 0.6681 & 0.4690 & 0.995 & 0.998 & 0.493 & 1.000 & 0.781 & 0.400 & 0.009 \\
{}BoostPFN TabPFN-v1 약학습기 T=\allowbreak{}50 × 500행 & 0.7368 & 0.5049 & 0.993 & 0.995 & 0.916 & 0.992 & 0.742 & 0.324 & 0.196 \\
{}LoCalPFN TabPFN-kNN: 행마다 국소 컨텍스트 1,000행 & 0.7121 & 0.5386 & 0.994 & 0.995 & 0.736 & 0.866 & 0.773 & 0.396 & 0.225 \\
\bottomrule
\end{longtable}
\endgroup
\end{landscape}
\begin{landscape}
\begingroup
\scriptsize
\begin{longtable}{>{\raggedright\arraybackslash}p{\dimexpr 0.23627\linewidth-2\tabcolsep\relax}>{\raggedright\arraybackslash}p{\dimexpr 0.08922\linewidth-2\tabcolsep\relax}>{\raggedright\arraybackslash}p{\dimexpr 0.08431\linewidth-2\tabcolsep\relax}>{\raggedright\arraybackslash}p{\dimexpr 0.08431\linewidth-2\tabcolsep\relax}>{\raggedright\arraybackslash}p{\dimexpr 0.08431\linewidth-2\tabcolsep\relax}>{\raggedright\arraybackslash}p{\dimexpr 0.08431\linewidth-2\tabcolsep\relax}>{\raggedright\arraybackslash}p{\dimexpr 0.08431\linewidth-2\tabcolsep\relax}>{\raggedright\arraybackslash}p{\dimexpr 0.08431\linewidth-2\tabcolsep\relax}>{\raggedright\arraybackslash}p{\dimexpr 0.08431\linewidth-2\tabcolsep\relax}>{\raggedright\arraybackslash}p{\dimexpr 0.08431\linewidth-2\tabcolsep\relax}}
\toprule
방법 & macro-F1 & tail F1 & benign & bot & brute & ddos & dos & inf & web \\ \midrule
\endfirsthead
방법 & macro-F1 & tail F1 & benign & bot & brute & ddos & dos & inf & web \\ \midrule
\endhead
{}다수 라벨 참조값 bank 없음 · test 벡터별 다수 라벨 직접 대입 & 1.0000 & 1.0000 & 1.000 & 1.000 & 1.000 & 1.000 & 1.000 & 1.000 & 1.000 \\
{}XGBoost train 풀 전체 (300 trees · depth 8) & 0.8467 & 0.6759 & 0.992 & 0.998 & 1.000 & 0.909 & 0.998 & 0.618 & 0.411 \\
{}XGBoost 같은 C0 100k행 (정보 예산 짝지음) & 0.7774 & 0.4854 & 0.994 & 0.997 & 1.000 & 1.000 & 0.992 & 0.382 & 0.077 \\
{}TabPFN-v3 무작위 100k 컨텍스트 (비율 보존) & 0.8055 & 0.5509 & 0.998 & 0.998 & 1.000 & 0.995 & 0.992 & 0.655 & 0.000 \\
{}ours 구성된 C0 100k (benign 0.75 · 공격 균등 · n\_\allowbreak{}est 4) & 0.8140 & 0.5702 & 0.998 & 0.998 & 1.000 & 0.998 & 0.992 & 0.619 & 0.094 \\
{}DistPFN ours(n\_\allowbreak{}est 4) 확률에 prior 보정 & 0.8106 & 0.5714 & 0.997 & 0.998 & 1.000 & 0.998 & 0.964 & 0.623 & 0.093 \\
{}BoostPFN TabPFN-v1 약학습기 T=\allowbreak{}50 × 500행 & 0.7708 & 0.4856 & 0.992 & 0.992 & 1.000 & 0.987 & 0.960 & 0.292 & 0.173 \\
\bottomrule
\end{longtable}
\endgroup
\end{landscape}
\begin{landscape}
\begingroup
\scriptsize
\begin{longtable}{>{\raggedright\arraybackslash}p{\dimexpr 0.23627\linewidth-2\tabcolsep\relax}>{\raggedright\arraybackslash}p{\dimexpr 0.08922\linewidth-2\tabcolsep\relax}>{\raggedright\arraybackslash}p{\dimexpr 0.08431\linewidth-2\tabcolsep\relax}>{\raggedright\arraybackslash}p{\dimexpr 0.08431\linewidth-2\tabcolsep\relax}>{\raggedright\arraybackslash}p{\dimexpr 0.08431\linewidth-2\tabcolsep\relax}>{\raggedright\arraybackslash}p{\dimexpr 0.08431\linewidth-2\tabcolsep\relax}>{\raggedright\arraybackslash}p{\dimexpr 0.08431\linewidth-2\tabcolsep\relax}>{\raggedright\arraybackslash}p{\dimexpr 0.08431\linewidth-2\tabcolsep\relax}>{\raggedright\arraybackslash}p{\dimexpr 0.08431\linewidth-2\tabcolsep\relax}>{\raggedright\arraybackslash}p{\dimexpr 0.08431\linewidth-2\tabcolsep\relax}}
\toprule
방법 & macro-F1 & tail F1 & benign & bot & brute & ddos & dos & inf & web \\ \midrule
\endfirsthead
방법 & macro-F1 & tail F1 & benign & bot & brute & ddos & dos & inf & web \\ \midrule
\endhead
{}다수 라벨 참조값 bank 없음 · test 벡터별 다수 라벨 직접 대입 & 1.0000 & 1.0000 & 1.000 & 1.000 & 1.000 & 1.000 & 1.000 & 1.000 & 1.000 \\
{}XGBoost train 풀 전체 (300 trees · depth 8) & 0.7968 & 0.5620 & 0.988 & 0.998 & 1.000 & 0.904 & 1.000 & 0.306 & 0.382 \\
{}XGBoost 같은 C0 100k행 (정보 예산 짝지음) & 0.7921 & 0.5177 & 0.993 & 0.997 & 0.999 & 0.999 & 1.000 & 0.169 & 0.387 \\
{}TabPFN-v3 무작위 100k 컨텍스트 (비율 보존) & 0.7590 & 0.4416 & 0.998 & 0.998 & 1.000 & 0.999 & 0.992 & 0.217 & 0.109 \\
{}ours 구성된 C0 100k (benign 0.75 · 공격 균등 · n\_\allowbreak{}est 4) & 0.7923 & 0.5257 & 0.996 & 0.998 & 1.000 & 0.982 & 0.992 & 0.205 & 0.374 \\
{}DistPFN ours(n\_\allowbreak{}est 4) 확률에 prior 보정 & 0.7931 & 0.5221 & 0.997 & 0.998 & 1.000 & 0.997 & 0.992 & 0.199 & 0.370 \\
{}BoostPFN TabPFN-v1 약학습기 T=\allowbreak{}50 × 500행 & 0.7194 & 0.3668 & 0.990 & 0.991 & 0.999 & 0.992 & 0.955 & 0.049 & 0.060 \\
\bottomrule
\end{longtable}
\endgroup
\end{landscape}
\begin{landscape}
\begingroup
\scriptsize
\begin{longtable}{>{\raggedright\arraybackslash}p{\dimexpr 0.23627\linewidth-2\tabcolsep\relax}>{\raggedright\arraybackslash}p{\dimexpr 0.08922\linewidth-2\tabcolsep\relax}>{\raggedright\arraybackslash}p{\dimexpr 0.08431\linewidth-2\tabcolsep\relax}>{\raggedright\arraybackslash}p{\dimexpr 0.08431\linewidth-2\tabcolsep\relax}>{\raggedright\arraybackslash}p{\dimexpr 0.08431\linewidth-2\tabcolsep\relax}>{\raggedright\arraybackslash}p{\dimexpr 0.08431\linewidth-2\tabcolsep\relax}>{\raggedright\arraybackslash}p{\dimexpr 0.08431\linewidth-2\tabcolsep\relax}>{\raggedright\arraybackslash}p{\dimexpr 0.08431\linewidth-2\tabcolsep\relax}>{\raggedright\arraybackslash}p{\dimexpr 0.08431\linewidth-2\tabcolsep\relax}>{\raggedright\arraybackslash}p{\dimexpr 0.08431\linewidth-2\tabcolsep\relax}}
\toprule
방법 & macro-F1 & tail F1 & benign & bot & brute & ddos & dos & inf & web \\ \midrule
\endfirsthead
방법 & macro-F1 & tail F1 & benign & bot & brute & ddos & dos & inf & web \\ \midrule
\endhead
{}다수 라벨 참조값 bank 없음 · test 벡터별 다수 라벨 직접 대입 & 0.9974 & 0.9939 & 1.000 & 1.000 & 1.000 & 1.000 & 1.000 & 0.989 & 0.992 \\
{}XGBoost train 풀 전체 (300 trees · depth 8) & 0.7613 & 0.4450 & 0.995 & 0.982 & 0.999 & 1.000 & 1.000 & 0.178 & 0.175 \\
{}XGBoost 같은 C0 100k행 (정보 예산 짝지음) & 0.8043 & 0.5482 & 0.989 & 0.970 & 0.999 & 1.000 & 0.997 & 0.131 & 0.543 \\
{}TabPFN-v3 무작위 100k 컨텍스트 (비율 보존) & 0.7411 & 0.3977 & 0.997 & 0.974 & 1.000 & 1.000 & 0.998 & 0.078 & 0.142 \\
{}ours 구성된 C0 100k (benign 0.75 · 공격 균등 · n\_\allowbreak{}est 4) & 0.8055 & 0.5678 & 0.993 & 0.985 & 1.000 & 0.945 & 0.997 & 0.249 & 0.469 \\
{}DistPFN ours(n\_\allowbreak{}est 4) 확률에 prior 보정 & 0.7590 & 0.4710 & 0.991 & 0.815 & 1.000 & 0.911 & 0.998 & 0.121 & 0.478 \\
{}BoostPFN TabPFN-v1 약학습기 T=\allowbreak{}50 × 500행 & 0.7404 & 0.4033 & 0.992 & 0.924 & 1.000 & 0.994 & 0.987 & 0.071 & 0.214 \\
{}LoCalPFN TabPFN-kNN: 행마다 국소 컨텍스트 1,000행 & 0.7690 & 0.5023 & 0.989 & 0.957 & 0.954 & 0.941 & 0.992 & 0.155 & 0.395 \\
\bottomrule
\end{longtable}
\endgroup
\end{landscape}
\begin{landscape}
\begingroup
\scriptsize
\begin{longtable}{>{\raggedright\arraybackslash}p{\dimexpr 0.23627\linewidth-2\tabcolsep\relax}>{\raggedright\arraybackslash}p{\dimexpr 0.08922\linewidth-2\tabcolsep\relax}>{\raggedright\arraybackslash}p{\dimexpr 0.08431\linewidth-2\tabcolsep\relax}>{\raggedright\arraybackslash}p{\dimexpr 0.08431\linewidth-2\tabcolsep\relax}>{\raggedright\arraybackslash}p{\dimexpr 0.08431\linewidth-2\tabcolsep\relax}>{\raggedright\arraybackslash}p{\dimexpr 0.08431\linewidth-2\tabcolsep\relax}>{\raggedright\arraybackslash}p{\dimexpr 0.08431\linewidth-2\tabcolsep\relax}>{\raggedright\arraybackslash}p{\dimexpr 0.08431\linewidth-2\tabcolsep\relax}>{\raggedright\arraybackslash}p{\dimexpr 0.08431\linewidth-2\tabcolsep\relax}>{\raggedright\arraybackslash}p{\dimexpr 0.08431\linewidth-2\tabcolsep\relax}}
\toprule
방법 & macro-F1 & tail F1 & benign & bot & brute & ddos & dos & inf & web \\ \midrule
\endfirsthead
방법 & macro-F1 & tail F1 & benign & bot & brute & ddos & dos & inf & web \\ \midrule
\endhead
{}다수 라벨 참조값 bank 없음 · test 벡터별 다수 라벨 직접 대입 & 1.0000 & 1.0000 & 1.000 & 1.000 & 1.000 & 1.000 & 1.000 & 1.000 & 1.000 \\
{}XGBoost train 풀 전체 (300 trees · depth 8) & 0.7759 & 0.4785 & 0.997 & 0.977 & 0.999 & 1.000 & 1.000 & 0.214 & 0.244 \\
{}XGBoost 같은 C0 100k행 (정보 예산 짝지음) & 0.7956 & 0.5278 & 0.990 & 0.973 & 0.999 & 1.000 & 0.997 & 0.100 & 0.511 \\
{}TabPFN-v3 무작위 100k 컨텍스트 (비율 보존) & 0.7321 & 0.3764 & 0.998 & 0.974 & 1.000 & 1.000 & 0.998 & 0.116 & 0.039 \\
{}ours 구성된 C0 100k (benign 0.75 · 공격 균등 · n\_\allowbreak{}est 4) & 0.7977 & 0.5344 & 0.996 & 0.982 & 1.000 & 0.987 & 0.998 & 0.125 & 0.496 \\
{}DistPFN ours(n\_\allowbreak{}est 4) 확률에 prior 보정 & 0.6150 & 0.2812 & 0.973 & 0.706 & 0.565 & 0.943 & 0.980 & 0.135 & 0.003 \\
{}BoostPFN TabPFN-v1 약학습기 T=\allowbreak{}50 × 500행 & 0.6884 & 0.2877 & 0.995 & 0.801 & 1.000 & 0.988 & 0.972 & 0.062 & 0.000 \\
{}LoCalPFN TabPFN-kNN: 행마다 국소 컨텍스트 1,000행 & 0.7588 & 0.4499 & 0.993 & 0.951 & 0.999 & 0.985 & 0.985 & 0.143 & 0.256 \\
\bottomrule
\end{longtable}
\endgroup
\end{landscape}
\begin{landscape}
\begingroup
\scriptsize
\begin{longtable}{>{\raggedright\arraybackslash}p{\dimexpr 0.23627\linewidth-2\tabcolsep\relax}>{\raggedright\arraybackslash}p{\dimexpr 0.08922\linewidth-2\tabcolsep\relax}>{\raggedright\arraybackslash}p{\dimexpr 0.08431\linewidth-2\tabcolsep\relax}>{\raggedright\arraybackslash}p{\dimexpr 0.08431\linewidth-2\tabcolsep\relax}>{\raggedright\arraybackslash}p{\dimexpr 0.08431\linewidth-2\tabcolsep\relax}>{\raggedright\arraybackslash}p{\dimexpr 0.08431\linewidth-2\tabcolsep\relax}>{\raggedright\arraybackslash}p{\dimexpr 0.08431\linewidth-2\tabcolsep\relax}>{\raggedright\arraybackslash}p{\dimexpr 0.08431\linewidth-2\tabcolsep\relax}>{\raggedright\arraybackslash}p{\dimexpr 0.08431\linewidth-2\tabcolsep\relax}>{\raggedright\arraybackslash}p{\dimexpr 0.08431\linewidth-2\tabcolsep\relax}}
\toprule
방법 & macro-F1 & tail F1 & benign & bot & brute & ddos & dos & inf & web \\ \midrule
\endfirsthead
방법 & macro-F1 & tail F1 & benign & bot & brute & ddos & dos & inf & web \\ \midrule
\endhead
{}다수 라벨 참조값 bank 없음 · test 벡터별 다수 라벨 직접 대입 & 1.0000 & 1.0000 & 1.000 & 1.000 & 1.000 & 1.000 & 1.000 & 1.000 & 1.000 \\
{}XGBoost train 풀 전체 (300 trees · depth 8) & 0.7763 & 0.4794 & 0.997 & 0.978 & 0.999 & 1.000 & 1.000 & 0.213 & 0.247 \\
{}XGBoost 같은 C0 100k행 (정보 예산 짝지음) & 0.7919 & 0.5191 & 0.990 & 0.973 & 0.997 & 1.000 & 0.998 & 0.096 & 0.488 \\
{}TabPFN-v3 무작위 100k 컨텍스트 (비율 보존) & 0.7541 & 0.4278 & 0.998 & 0.972 & 1.000 & 1.000 & 0.998 & 0.152 & 0.160 \\
{}ours 구성된 C0 100k (benign 0.75 · 공격 균등 · n\_\allowbreak{}est 4) & 0.8039 & 0.5460 & 0.995 & 0.983 & 1.000 & 0.996 & 0.998 & 0.178 & 0.477 \\
{}DistPFN ours(n\_\allowbreak{}est 4) 확률에 prior 보정 & 0.6041 & 0.2715 & 0.976 & 0.625 & 0.468 & 0.973 & 0.996 & 0.186 & 0.003 \\
{}BoostPFN TabPFN-v1 약학습기 T=\allowbreak{}50 × 500행 & 0.6636 & 0.2772 & 0.986 & 0.792 & 1.000 & 0.863 & 0.966 & 0.040 & 0.000 \\
{}LoCalPFN TabPFN-kNN: 행마다 국소 컨텍스트 1,000행 & 0.7402 & 0.4527 & 0.985 & 0.952 & 0.999 & 0.862 & 0.977 & 0.122 & 0.285 \\
\bottomrule
\end{longtable}
\endgroup
\end{landscape}
\begin{landscape}
\begingroup
\scriptsize
\begin{longtable}{>{\raggedright\arraybackslash}p{\dimexpr 0.15293\linewidth-2\tabcolsep\relax}>{\raggedright\arraybackslash}p{\dimexpr 0.12633\linewidth-2\tabcolsep\relax}>{\raggedright\arraybackslash}p{\dimexpr 0.09973\linewidth-2\tabcolsep\relax}>{\raggedright\arraybackslash}p{\dimexpr 0.09973\linewidth-2\tabcolsep\relax}>{\raggedright\arraybackslash}p{\dimexpr 0.09973\linewidth-2\tabcolsep\relax}>{\raggedright\arraybackslash}p{\dimexpr 0.09973\linewidth-2\tabcolsep\relax}>{\raggedright\arraybackslash}p{\dimexpr 0.11037\linewidth-2\tabcolsep\relax}>{\raggedright\arraybackslash}p{\dimexpr 0.21144\linewidth-2\tabcolsep\relax}}
\toprule
데이터 & ours seed 42 & seed 43 & seed 44 & 평균 & 폭 & XGB-c0 평균 & ours − XGB-c0 (42 /\allowbreak{} 43 /\allowbreak{} 44) \\ \midrule
\endfirsthead
데이터 & ours seed 42 & seed 43 & seed 44 & 평균 & 폭 & XGB-c0 평균 & ours − XGB-c0 (42 /\allowbreak{} 43 /\allowbreak{} 44) \\ \midrule
\endhead
{}원본 & 0.7631 & 0.7621 & 0.7609 & 0.7620 & 0.0023 & 0.7498 & +\allowbreak{}0.0107 /\allowbreak{} +\allowbreak{}0.0134 /\allowbreak{} +\allowbreak{}0.0126 \\
{}중복 제거 & 0.8062 & 0.7985 & 0.8056 & 0.8035 & 0.0077 & 0.8110 & +\allowbreak{}0.0019 /\allowbreak{} -0.0175 /\allowbreak{} -0.0068 \\
{}중복 제거 +\allowbreak{} 모순은 다수 라벨 & 0.7994 & 0.8054 & 0.7858 & 0.7969 & 0.0196 & 0.7965 & +\allowbreak{}0.0038 /\allowbreak{} +\allowbreak{}0.0093 /\allowbreak{} -0.0118 \\
{}중복 제거 +\allowbreak{} 모순 벡터 삭제 & 0.7975 & 0.7783 & 0.7689 & 0.7816 & 0.0286 & 0.7916 & +\allowbreak{}0.0056 /\allowbreak{} -0.0123 /\allowbreak{} -0.0233 \\
\bottomrule
\end{longtable}
\endgroup
\end{landscape}
100k C0는 seed에 따라 macro가 0.675–0.765를 오간 기록이 있어(0902 §8ag) 처리별 차이를 읽기 전에 seed 폭을 먼저 본다. 위 방법 비교표는 seed 42 한 번이므로, 처리 간 차이가 이 폭보다 작으면 차이로 읽지 않는다. 같은 C0를 준 XGBoost와 비교하면 ours가 세 seed 모두 앞서는 arm은 원본뿐이고, 중복 제거 세 arm에서는 seed 42만 ours가 앞선다 — 위 표의 "ours 4/\allowbreak{}4 최고"는 seed 42 한 번의 결과이며, dedup arm에서 XGB-c0 대비 우위는 seed 폭 안이라 차이로 읽지 않는다(0915 재검). 중복 유지 +\allowbreak{} 모순 제거 두 arm(EXP44)은 seed 42 단일이라 폭이 없다.

여기서 ours는 구성된 C0(100k · benign 0.75 · 공격 균등 · n\_\allowbreak{}estimators 4) 위의 frozen TabPFN-v3 이고 그 이상이 아니다 — scorer·verifier 층은 EXP38·39에서 6조건 모두 global을 선택했다(Δglobal =\allowbreak{} 0). 따라서 이 표는 C0 구성의 비교다. 라우팅 개선 가능성은 담당 배정·실제 S/\allowbreak{}V·bank 진단을 분리한 0918 수정 보고서에서 판단한다.

EXP43 (설계, 실행 보류). tabpfn/\allowbreak{}scripts/\allowbreak{}nfv3\_\allowbreak{}v3\_\allowbreak{}exp43\_\allowbreak{}ours\_\allowbreak{}majfill.py 의 knob --attack-fill 은 예산·클래스 균형·benign 행을 고정한 채 공격 클래스 할당량을 어떤 행으로 채울지 만 바꾼다( majority\_\allowbreak{}consistent : 그 벡터의 다수 라벨이 해당 클래스인 행 우선, clean : 다른 라벨이 전혀 없는 벡터 우선). 근거는 §1의 감사다 — infiltration 행의 75.1\%, web의 69.6\%가 상반 라벨 트윈이라 균일 추첨은 infiltration 슬롯의 3/\allowbreak{}4을 benign과 구분 불가능한 벡터에 쓴다. 단 현재 구현은 다수 라벨을 train+\allowbreak{}val+\allowbreak{}test 전체에서 계산하므로 그대로는 selective snooping이다 — 실행한다면 train 풀 다수 라벨 로 바꾼 새 frozen 스크립트여야 한다. 2026-09-14 사용자 판단으로 실행은 보류했다.

대조군 검증에서 나온 것 (실행함). --attack-fill natural 로 중복 제거 arm을 돌려(run 20260914\_\allowbreak{}154808 ) EXP42의 ours와 같은 행을 뽑는지 확인했다. 클래스별 목표 수(benign 75,000 /\allowbreak{} bot 841 /\allowbreak{} brute 5,967 /\allowbreak{} ddos 5,966 /\allowbreak{} dos 5,966 /\allowbreak{} inf 5,965 /\allowbreak{} web 295)와 추첨 seed가 같아 컨텍스트 행은 동일한데 macro-F1은 0.8055(EXP42) vs 0.8087(EXP43)로 0.0032 어긋난다. 원인은 TabPFN config 추첨의 전역 RNG 상태다(EXP42는 같은 프로세스에서 tabpfn\_\allowbreak{}plain이 먼저 돌았다). 함의: 컨텍스트 knob 비교는 반드시 같은 스크립트·같은 프로세스 순서의 natural 대조군과 해야 한다. EXP42 값과 비교하면 ±0.003의 가짜 개선을 볼 수 있다.

방향 (09-15). 데이터 감사 실험은 EXP44로 마치고 ours 자체의 개선 실험으로 돌아간다. 후보와 우선순위는 docs/\allowbreak{}lablog/\allowbreak{}reports/\allowbreak{}20260915.tex 에 기록하며, 실행은 사용자와 합의한 뒤 시작한다.

\subsection*{6. 시간순 클래스 구성}
CIC-IDS2018·ToN-IoT·BoT-IoT는 공격 family가 날짜 블록 으로 실행됐다(arXiv 2503.04404 Table 4). benign만 전 구간에 걸쳐 있고 각 공격은 자기 날짜 안에 있다. 파이프라인의 family별 시간순 split은 family 내부의 앞뒤를 자른 것이고, 새 family가 나중에 등장하는 상황은 family 등장 순서로 따로 구성해야 한다(ToN-IoT: scanning  \ensuremath{\rightarrow}  …  \ensuremath{\rightarrow}  backdoor/\allowbreak{}ransomware  \ensuremath{\rightarrow}  mitm). family 내부에도 변화는 있다: BoT-IoT ddos의 test 벡터는 train에 사실상 없고, ToN-IoT scanning·BoT-IoT theft는 같은 벡터의 라벨이 시점에 따라 바뀐다. UNSW-NB15만 모든 family가 두 수집일에 동시에 존재한다.

\begingroup
\footnotesize
\begin{longtable}{>{\raggedright\arraybackslash}p{\dimexpr 0.16900\linewidth-2\tabcolsep\relax}>{\raggedright\arraybackslash}p{\dimexpr 0.27700\linewidth-2\tabcolsep\relax}>{\raggedright\arraybackslash}p{\dimexpr 0.27700\linewidth-2\tabcolsep\relax}>{\raggedright\arraybackslash}p{\dimexpr 0.27700\linewidth-2\tabcolsep\relax}}
\toprule
데이터셋 & 결함의 본질 & 정답 참조 진단에서 구분할 점 & 쓸 수 있는 질문 \\ \midrule
\endfirsthead
데이터셋 & 결함의 본질 & 정답 참조 진단에서 구분할 점 & 쓸 수 있는 질문 \\ \midrule
\endhead
{}CIC-IDS2018 & 시간 창 라벨링  \ensuremath{\rightarrow}  infiltration·web이 benign과 동일 벡터; slowhttptest=\allowbreak{}ftp\_\allowbreak{}bruteforce 단일 벡터 & infiltration 0.97 vs 다수 라벨 참조 0.67, web 0.28 vs 0.27 & context 구성(bot·ddos·brute의 깨끗한 부분); 새 family 등장은 bot(고유 벡터 100\%, 마지막 날); infiltration은 3분해 보고와 함께만 \\
{}ToN-IoT & IP 대역+\allowbreak{}시간 창 라벨링  \ensuremath{\rightarrow}  benign 62\%·ddos 78\%·scanning 97\%가 충돌; 응답 없는 SYN 한 개가 수백만 benign 행 & scanning은 test 안에서는 순수(다수 라벨 참조 1.00)하나 train 라벨과 98\% 충돌  \ensuremath{\rightarrow}  시간 이동 진단 필요 & mitm·ransomware 정밀도(다수 라벨 참조값 0.99/\allowbreak{}0.95 vs XGB 0.16/\allowbreak{}0.15); 늦게 등장하는 family +\allowbreak{} k 라벨 refresh \\
{}BoT-IoT & 합성 volumetric, 중복률 96\%; benign 0.3\%가 recon/\allowbreak{}dos와 충돌; theft 벡터의 라벨이 5월 recon \ensuremath{\rightarrow} 6월 theft & theft 다수 라벨 참조 0.98지만 train 라벨과 75\% 충돌; benign 0.85 & sanity check 이상 없음 \\
{}UNSW-NB15 & TTL 지름길로 benign/\allowbreak{}attack 자명; 공격 카테고리는 IXIA 익스플로잇 목록 라벨이라 서로 근사 동일 & backdoor·dos·analysis의 카테고리 혼동(근사 중복이라 oracle 아래 여지가 학습 가능하다는 보장이 없음) & 공격 간 꼬리(보조, 주의 문구 필수) \\
\bottomrule
\end{longtable}
\endgroup
답은 데이터 교체가 아니라 평가를 바꾸는 것이다. 이는 새 프로토콜이 아니라 이미 요구되는 관행이다(Arp et al. USENIX Sec'22의 temporal snooping·spurious correlation, Flood et al. EuroS\&P'24의 near-duplicate·traffic collapse·wrong label 진단, Liu et al. CNS'22의 라벨 감사, Allamanis Onward!'19의 중복/\allowbreak{}비중복 test 분해, El Mahdaouy et al. arXiv 2602.05594). 구체적으로: (a) headline macro-F1은 전 클래스로 보고하되 클래스별 다수 라벨 참조값과 train 라벨 충돌률을 옆에 둔다(꼬리 클래스를 빼지 않는다); (b) test를 "train에 같은 라벨로 있던 벡터 /\allowbreak{} 다른 라벨로 있던 벡터 /\allowbreak{} 처음 보는 벡터"로 나눠 보고한다; (c) split은 시간순이며 family 등장 순서 프로토콜을 구분한다; (d) 충돌 벡터는 행 단위 학습이 아니라 그룹 단위 비용 정책으로 다루되 다음 시간 구간에서 forward로 검증한다; (e) 다수 라벨 참조값과 bank 정답 포함률은 진단으로 쓰며 그 격차를 학습 가능한 headroom으로 주장하지 않는다; (f) 대상 클래스는 ToN mitm·ransomware, CIC web(꼬리 정밀도), ToN d9–d10과 CIC bot(새 family 등장)으로 미리 고른다.

모순 처리에 대한 판정(09-15). 모순은 동일 벡터에 같은 답을 내는 분류기의 정확일치 제약이지만, 지우거나(EXP41 중복 제거·벡터 삭제, EXP44 벡터 삭제) 고치는(EXP41·EXP44 다수 라벨) 어느 처리도 꼬리를 지키지 못했고, 처리 뒤에는 방법 순위까지 뒤집힌다(§4). 다수 라벨 재라벨은 test를 들여다보는 selective snooping이라 배포 전처리가 될 수 없고, 문헌도 IDS 벤치마크의 모순을 지우지 않고 원본 위에서 보고하는 것이 관행이다. 따라서 headline은 원본 벤치마크 로 두고, 모순은 전처리가 아니라 (a)·(b)의 보고 축 — 클래스별 다수 라벨 참조값과 test 3분해 — 로 다룬다. 처리된 벤치마크(§4)는 보충 자료로만 남긴다. 모순에서 얻은 지식은 train 풀 쪽에서만 쓴다(예: 컨텍스트 구성 시 train 다수 라벨이 그 클래스인 행을 우선).

입력 data/\allowbreak{}nfv3\_\allowbreak{}energy\_\allowbreak{}suite\_\allowbreak{}uncapped\_\allowbreak{}scenarios.pkl . hash =\allowbreak{} pd.util.hash\_\allowbreak{}pandas\_\allowbreak{}object (nan\_\allowbreak{}to\_\allowbreak{}num float32 행). 감사 코드와 CSV: docs/\allowbreak{}research/\allowbreak{}20260911/\allowbreak{}dataset\_\allowbreak{}quality\_\allowbreak{}audit/\allowbreak{} (audit\_\allowbreak{}pass1.py, audit\_\allowbreak{}pass2\_\allowbreak{}ceilings\_\allowbreak{}unsw\_\allowbreak{}xgb.py, unsw\_\allowbreak{}probe\_\allowbreak{}ttl\_\allowbreak{}confusion.py, pair\_\allowbreak{}conflicts.py  \ensuremath{\rightarrow}  pair\_\allowbreak{}conflicts.json, ceilings\_\allowbreak{}v2.py  \ensuremath{\rightarrow}  ceilings\_\allowbreak{}v2.csv(다수 라벨 참조값 =\allowbreak{} ceil\_\allowbreak{}majority\_\allowbreak{}test), cic\_\allowbreak{}bank\_\allowbreak{}oracle.py  \ensuremath{\rightarrow}  cic\_\allowbreak{}bank\_\allowbreak{}oracle\_\allowbreak{}derivation.json, report\_\allowbreak{}data\_\allowbreak{}twins/\allowbreak{}time.py, build\_\allowbreak{}dataset\_\allowbreak{}report.py). EXP41: scripts/\allowbreak{}exp41\_\allowbreak{}dedup\_\allowbreak{}dataset\_\allowbreak{}xgb.py (results/\allowbreak{}…\_\allowbreak{}exp41\_\allowbreak{}dedup\_\allowbreak{}dataset\_\allowbreak{}xgb) · tabpfn/\allowbreak{}scripts/\allowbreak{}nfv3\_\allowbreak{}v3\_\allowbreak{}exp41\_\allowbreak{}dedup\_\allowbreak{}dataset\_\allowbreak{}tabpfn.py (4 처리 × seed 42·43·44). EXP42(방법 비교): tabpfn/\allowbreak{}scripts/\allowbreak{}nfv3\_\allowbreak{}v3\_\allowbreak{}exp42\_\allowbreak{}dedup\_\allowbreak{}methods.py — EXP41과 같은 dedup·재split 코드를 그대로 써서 arm별 train/\allowbreak{}val/\allowbreak{}test 행이 동일하고, XGBoost train 전체는 EXP41 run의 예측을 test\_\allowbreak{}idx 일치 확인 후 재사용한다. SOTA는 tabpfn/\allowbreak{}third\_\allowbreak{}party/\allowbreak{}\{BoostPFN,LoCalPFN\} 의 TabPFN-v1을 쓰며 exp35·exp36과 같은 설정이다. 달성 F1 출처: CIC-IDS2018은 EXP39 test global(0909 §3-C), BoT/\allowbreak{}ToN은 0818 full-pool XGBoost run, UNSW는 0911 XGB. EXP43(ours v2 설계 · natural 대조군만 실행): tabpfn/\allowbreak{}scripts/\allowbreak{}nfv3\_\allowbreak{}v3\_\allowbreak{}exp43\_\allowbreak{}ours\_\allowbreak{}majfill.py . EXP44(중복 유지 +\allowbreak{} 모순 제거): tabpfn/\allowbreak{}scripts/\allowbreak{}nfv3\_\allowbreak{}v3\_\allowbreak{}exp44\_\allowbreak{}conflict\_\allowbreak{}only.py (relabel\_\allowbreak{}majority · drop\_\allowbreak{}vectors, LoCalPFN 제외). 랩 로그 docs/\allowbreak{}lablog/\allowbreak{}reports/\allowbreak{}20260911.tex · 0914.md · 0915.md .



\section*{출처와 실행 추적}

상세 수치의 원시 실행은 \path{inventory/runs.csv}, 실험별 연결은 \path{inventory/experiments.csv}에서 확인한다. 표·그림의 재현 입력은 \path{evidence/}와 \path{data/}에 둔다. 문서 이동·통합 내역은 \path{configs/experiments/migrations/20261020_documents/manifest.json}에 기록했다.

\end{document}
